





\documentclass[11pt]{article}
\usepackage[T1]{fontenc}
\usepackage[utf8]{inputenc}
\usepackage{lmodern}
\usepackage[english]{babel}
\usepackage{amsmath,amssymb,mathtools,mathrsfs,array,booktabs,pdflscape,adjustbox,diagbox,graphicx}
\usepackage{geometry}
\usepackage{microtype}
\usepackage{newunicodechar}
\newunicodechar{—}{---}
\newunicodechar{“}{``}
\newunicodechar{”}{''}
\usepackage[hidelinks,unicode=true,pdfencoding=auto,bookmarksopen=true,bookmarksnumbered=false]{hyperref}
\geometry{a4paper,margin=1.45cm}
\setlength{\parindent}{1.25em}
\setlength{\parskip}{0pt}
\makeatletter
\renewcommand\footnotesize{\@setfontsize\footnotesize{7.5}{8.5}\selectfont}
\makeatother
\providecommand{\srcfn}[2]{%
  \begingroup
  \renewcommand{\thefootnote}{#1}%
  \footnote{#2}%
  \addtocounter{footnote}{-1}%
  \endgroup}
\providecommand{\srcfnmark}[1]{%
  \begingroup
  \renewcommand{\thefootnote}{#1}%
  \footnotemark%
  \addtocounter{footnote}{-1}%
  \endgroup}
\providecommand{\srcfntext}[2]{%
  \begingroup
  \renewcommand{\thefootnote}{#1}%
  \footnotetext{#2}%
  \endgroup}
\providecommand{\noethpIIrosette}{\makebox[9.44444pt][c]{\raisebox{-3.85pt}[0pt][0pt]{\includegraphics[width=14.726716981132075pt]{authority_rosette_native_supported_mask.png}}}}
\providecommand{\srcnumdisplay}[3][3.4em]{%
  \par\smallskip\noindent
  \begin{minipage}[t]{#1}\vspace*{0pt}#2\end{minipage}%
  \begin{minipage}[t]{\dimexpr\linewidth-#1\relax}\vspace*{0pt}\centering
  \(\displaystyle #3\)
  \end{minipage}\par\smallskip}
\providecommand{\qtext}[1]{``#1''}
\providecommand{\widebar}[1]{\overline{#1}}
\providecommand{\tightlist}{\setlength{\itemsep}{0pt}\setlength{\parskip}{0pt}}
\providecommand{\A}{\Delta}
\providecommand{\D}{\Delta}
\allowdisplaybreaks
\setlength{\emergencystretch}{10em}
\sloppy
\hbadness=10000
\vbadness=10000
\hfuzz=2pt
\vfuzz=10pt
\raggedbottom
\providecommand{\R}{\varrho}
\providecommand{\hata}{\widehat{a}}
\providecommand{\hatv}{\widehat{v}}
\providecommand{\modu}[1]{\;\operatorname{mod} #1}
\providecommand{\mcell}[1]{\ensuremath{\begin{gathered}#1\end{gathered}}}
\providecommand{\pair}[2]{(#1\,|\,#2)}
\providecommand{\eps}{\varepsilon}
\providecommand{\rhoidx}{\varrho}

\providecommand{\dd}{\mathrm d}
\providecommand{\Div}{\operatorname{Div}}
\providecommand{\G}{\mathfrak G}
\providecommand{\bd}{\bar{\delta}}
\providecommand{\p}{\partial}

\providecommand{\frS}{\mathfrak S}
\providecommand{\frp}{\mathfrak p}
\providecommand{\fra}{\mathfrak a}
\providecommand{\frb}{\mathfrak b}
\providecommand{\frc}{\mathfrak c}
\providecommand{\frf}{\mathfrak f}
\providecommand{\frr}{\mathfrak r}
\providecommand{\frm}{\mathfrak m}
\providecommand{\frD}{\mathfrak D}
\providecommand{\calA}{\mathcal A}
\providecommand{\calB}{\mathcal B}
\providecommand{\calN}{\mathcal N}


\providecommand{\p}{\partial}
\providecommand{\dd}{\mathrm d}
\providecommand{\modu}[1]{\;\operatorname{mod} #1}
\providecommand{\mM}{\mathfrak{M}}
\providecommand{\mN}{\mathfrak{N}}
\providecommand{\mL}{\mathfrak{L}}
\providecommand{\mG}{\mathfrak{G}}
\providecommand{\mH}{\mathfrak{H}}
\providecommand{\mH}{\mathfrak{H}}
\providecommand{\mA}{\mathfrak{A}}
\providecommand{\mB}{\mathfrak{B}}
\providecommand{\mX}{\mathfrak{X}}
\providecommand{\mY}{\mathfrak{Y}}
\providecommand{\ma}{\mathfrak{a}}
\providecommand{\mb}{\mathfrak{b}}
\providecommand{\mc}{\mathfrak{c}}
\providecommand{\mx}{\mathfrak{x}}
\providecommand{\my}{\mathfrak{y}}
\providecommand{\mz}{\mathfrak{z}}
\providecommand{\me}{\mathfrak{e}}
\providecommand{\mbarA}{\overline{\mathfrak{A}}}
\providecommand{\mbarB}{\overline{\mathfrak{B}}}
\providecommand{\mbar}[1]{\overline{#1}}
\providecommand{\mC}{\mathfrak{C}}
\providecommand{\mE}{\mathfrak{E}}
\providecommand{\mF}{\mathfrak{F}}
\providecommand{\mT}{\mathfrak{T}}
\providecommand{\mW}{\mathfrak{W}}
\providecommand{\md}{\mathfrak{d}}
\providecommand{\mf}{\mathfrak{f}}
\providecommand{\mm}{\mathfrak{m}}
\providecommand{\mn}{\mathfrak{n}}
\providecommand{\mo}{\mathfrak{o}}
\providecommand{\mpideal}{\mathfrak{p}}
\providecommand{\mq}{\mathfrak{q}}
\providecommand{\mt}{\mathfrak{t}}
\providecommand{\mv}{\mathfrak{v}}
\providecommand{\bara}{\overline{\mathfrak{a}}}

\providecommand{\mU}{\mathfrak{U}}
\providecommand{\mP}{\mathfrak{P}}
\providecommand{\mQ}{\mathfrak{Q}}
\providecommand{\mS}{\mathfrak{S}}
\providecommand{\mD}{\mathfrak{D}}
\providecommand{\mR}{\mathfrak{R}}
\providecommand{\mV}{\mathfrak{V}}
\providecommand{\ideal}[1]{\mathfrak{#1}}

\providecommand{\Amod}{\mathfrak A}
\providecommand{\Bmod}{\mathfrak B}
\providecommand{\Gmod}{\mathfrak G}
\providecommand{\Rmod}{\mathfrak R}
\providecommand{\Smod}{\mathfrak S}
\providecommand{\mideal}{\mathfrak m}
\providecommand{\nideal}{\mathfrak n}
\providecommand{\gideal}{\mathfrak g}
\providecommand{\bb}{\mathfrak b}
\providecommand{\cc}{\mathfrak c}
\providecommand{\zero}[1]{\equiv 0(#1)}
\providecommand{\Norm}{N}


\providecommand{\Amod}{\mathfrak A}
\providecommand{\Bmod}{\mathfrak B}
\providecommand{\Gmod}{\mathfrak G}
\providecommand{\Mmod}{\mathfrak M}
\providecommand{\Nmod}{\mathfrak N}
\providecommand{\Smod}{\mathfrak S}
\providecommand{\mideal}{\mathfrak m}
\providecommand{\nideal}{\mathfrak n}
\providecommand{\gideal}{\mathfrak g}
\providecommand{\hideal}{\mathfrak h}
\providecommand{\jideal}{\mathfrak j}
\providecommand{\tideal}{\mathfrak t}
\providecommand{\aideal}{\mathfrak a}
\providecommand{\bideal}{\mathfrak b}
\providecommand{\zero}[1]{\equiv 0\pmod{#1}}
\providecommand{\Norm}{N}


\providecommand{\frakS}{\mathfrak S}
\providecommand{\frakM}{\mathfrak M}
\providecommand{\frakG}{\mathfrak G}
\providecommand{\frakm}{\mathfrak m}
\providecommand{\frakn}{\mathfrak n}
\providecommand{\frakp}{\mathfrak p}
\providecommand{\frakq}{\mathfrak q}
\providecommand{\frakr}{\mathfrak r}
\providecommand{\fraka}{\mathfrak a}
\providecommand{\frakb}{\mathfrak b}
\providecommand{\frakc}{\mathfrak c}
\providecommand{\frakd}{\mathfrak d}
\providecommand{\frake}{\mathfrak e}
\providecommand{\frakg}{\mathfrak g}
\providecommand{\fraks}{\mathfrak s}
\providecommand{\frako}{\mathfrak o}
\providecommand{\frK}{\mathfrak K}
\providecommand{\barfrakm}{\overline{\mathfrak m}}

\providecommand{\frakh}{\mathfrak h}
\providecommand{\frakt}{\mathfrak t}

\providecommand{\frakR}{\mathfrak R}
\providecommand{\frakT}{\mathfrak T}
\providecommand{\frakH}{\mathfrak H}
\providecommand{\frakP}{\mathfrak P}
\providecommand{\frakQ}{\mathfrak Q}
\providecommand{\frakZ}{\mathfrak Z}
\providecommand{\frakN}{\mathfrak N}
\providecommand{\frakB}{\mathfrak B}
\providecommand{\frakF}{\mathfrak F}
\providecommand{\frakL}{\mathfrak L}
\providecommand{\frakW}{\mathfrak W}
\providecommand{\frakGg}{\mathfrak G}
\providecommand{\Endl}{\operatorname{End}}
\providecommand{\Gal}{\operatorname{Gal}}
\providecommand{\Trdeg}{\operatorname{trdeg}}

\providecommand{\frR}{\mathfrak R}
\providecommand{\frK}{\mathfrak K}
\providecommand{\frL}{\mathfrak L}
\providecommand{\frM}{\mathfrak M}
\providecommand{\frF}{\mathfrak F}
\providecommand{\frT}{\mathfrak T}
\providecommand{\frN}{\mathfrak N}
\providecommand{\frO}{\mathfrak O}
\providecommand{\frB}{\mathfrak B}
\providecommand{\frC}{\mathfrak C}
\providecommand{\frA}{\mathfrak A}

\providecommand{\mR}{\mathfrak{R}}
\providecommand{\bR}{\overline{\mathfrak{R}}}
\providecommand{\mK}{\mathfrak{K}}
\providecommand{\mL}{\mathfrak{L}}
\providecommand{\mS}{\mathfrak{S}}
\providecommand{\mT}{\mathfrak{T}}
\providecommand{\mZ}{\mathfrak{Z}}
\providecommand{\mO}{\mathfrak{O}}
\providecommand{\mo}{\mathfrak{o}}
\providecommand{\mD}{\mathfrak{D}}
\providecommand{\mE}{\mathfrak{E}}
\providecommand{\mP}{\mathfrak{P}}
\providecommand{\mf}{\mathfrak{f}}
\providecommand{\mpideal}{\mathfrak{p}}
\providecommand{\mq}{\mathfrak{q}}
\providecommand{\mt}{\mathfrak{t}}
\providecommand{\mOmega}{\mathfrak{\Omega}}
\providecommand{\bP}{\overline{P}}
\providecommand{\bq}{\overline{\mathfrak{q}}}
\providecommand{\bp}{\overline{\mathfrak{p}}}
\providecommand{\ba}{\overline{\mathfrak{a}}}
\providecommand{\bb}{\overline{\mathfrak{b}}}
\providecommand{\tr}{\operatorname{Tr}}
\providecommand{\Norm}{\operatorname{N}}


\providecommand{\rank}{\operatorname{rank}}

\providecommand{\frakA}{\mathfrak A}
\providecommand{\frakB}{\mathfrak B}
\providecommand{\frakM}{\mathfrak M}
\providecommand{\frakN}{\mathfrak N}
\providecommand{\frako}{\mathfrak o}
\providecommand{\fraka}{\mathfrak a}
\providecommand{\frakc}{\mathfrak c}
\providecommand{\frakp}{\mathfrak p}
\providecommand{\frakO}{\mathfrak O}
\providecommand{\frakP}{\mathfrak P}
\providecommand{\frakD}{\mathfrak D}
\providecommand{\frakG}{\mathfrak G}
\providecommand{\frakK}{\mathfrak K}
\providecommand{\Sp}{\operatorname{Sp}}
\providecommand{\rank}{\operatorname{rank}}
\providecommand{\Trdeg}{\operatorname{Trdeg}}
\providecommand{\charac}{\operatorname{char}}

\providecommand{\Om}{\Omega}
\providecommand{\Omfrakp}{\Omega_{\mathfrak p}}
\providecommand{\frakA}{\mathfrak A}
\providecommand{\frakB}{\mathfrak B}
\providecommand{\frakC}{\mathfrak C}
\providecommand{\frakG}{\mathfrak G}
\providecommand{\frakH}{\mathfrak H}
\providecommand{\frakP}{\mathfrak P}
\providecommand{\frakp}{\mathfrak p}
\providecommand{\frakq}{\mathfrak q}
\providecommand{\frakr}{\mathfrak r}
\providecommand{\frakS}{\mathfrak S}
\providecommand{\Cl}{\operatorname{Cl}}
\providecommand{\ind}{\operatorname{ind}}
\providecommand{\N}{\operatorname{N}}
\providecommand{\Mat}{\operatorname{M}}
\providecommand{\Gal}{\operatorname{Gal}}
\providecommand{\Norm}{\operatorname{N}}
\providecommand{\nr}{\operatorname{nr}}
\providecommand{\Tr}{\operatorname{Tr}}
\providecommand{\Br}{\operatorname{Br}}

\providecommand{\frR}{\mathfrak R}
\providecommand{\frS}{\mathfrak S}
\providecommand{\frT}{\mathfrak T}
\providecommand{\frM}{\mathfrak M}
\providecommand{\frN}{\mathfrak N}
\providecommand{\frA}{\mathfrak A}
\providecommand{\frB}{\mathfrak B}
\providecommand{\frX}{\mathfrak X}
\providecommand{\frY}{\mathfrak Y}
\providecommand{\frZ}{\mathfrak Z}
\providecommand{\frG}{\mathfrak G}
\providecommand{\frH}{\mathfrak H}
\providecommand{\frL}{\mathfrak L}
\providecommand{\frU}{\mathfrak U}
\providecommand{\frD}{\mathfrak D}
\providecommand{\frC}{\mathfrak C}
\providecommand{\frP}{\mathfrak P}
\providecommand{\Gcal}{\mathfrak G}
\providecommand{\Sp}{\operatorname{Sp}}

\providecommand{\Hh}{\mathfrak H}
\providecommand{\Ii}{\mathfrak I}
\providecommand{\Jj}{\mathfrak J}
\providecommand{\Aa}{\mathfrak A}
\providecommand{\Bb}{\mathfrak B}
\providecommand{\Cc}{\mathfrak C}
\providecommand{\Oo}{\mathfrak O}
\providecommand{\oo}{\mathfrak o}
\providecommand{\pp}{\mathfrak p}
\providecommand{\PP}{\mathfrak P}
\providecommand{\LL}{\mathfrak L}
\providecommand{\RR}{\mathfrak R}
\providecommand{\ZZ}{\mathfrak Z}
\providecommand{\ee}{\mathfrak e}

\providecommand{\Gg}{\mathfrak G}
\hypersetup{%
  pdftitle={Emmy Noether: Collected Mathematical Works in English},%
  pdfauthor={Emmy Noether},%
  pdfsubject={Complete maintained English corpus edition; source authority NOETH-DE-ED-0014},%
  pdfkeywords={Emmy Noether, English translation, collected works, algebra, invariant theory}%
}
\setcounter{tocdepth}{1}
\providecommand{\editionentry}[2]{%
  \phantomsection\addcontentsline{toc}{section}{#1}\label{#2}}
\providecommand{\editionpartentry}[2]{%
  \phantomsection\addcontentsline{toc}{subsection}{#1}\label{#2}}

\begin{document}

\section*{Elimination Theory and General Ideal Theory}
\begin{center}
Math. Ann. 90 (1923), pp. 229--261\\[0.9em]
{\Large\bfseries Elimination Theory and General Ideal Theory.}\\[1.0em]
By\\[0.45em]
Emmy Noether in Göttingen.\\[0.6em]
\rule{4em}{0.4pt}
\end{center}

In what follows the question is the placement of elimination theory -- in the arithmetical form that I gave to Hentzelt's presentation\footnote{Kurt Hentzelt, Zur Theorie der Polynomideale und Resultanten. Edited by E. Noether. Math. Ann. 88 (1922), pp. 53--79; cited H.-N.} and that may be described as an ideal theory in the polynomial domain -- within general ideal theory,\footnote{E. Noether, Idealtheorie in Ringbereichen, Math. Ann. 83 (1921), pp. 24--66; cited Ideal Theory.} and, at the same time, a new foundation of the parts referring to zeroes. The basic concepts of both theories are briefly assembled in § 1, as are those of field theory,\footnote{E. Steinitz, Algebraische Theorie der Körper, J. f. M. 137 (1910), pp. 167--309; cited Steinitz.} so that I can refer to them here.

The classification and the new foundation group themselves around four questions: the arithmetical formulation of the concept of dimension; the theory of zeroes; the connection between the decomposition theorems for norms and elementary divisors and those of general ideal theory; the characterization of prime ideals and primary ideals by norm and elementary-divisor form.

The arithmetical formulation of the \emph{concept of dimension} (§ 4) is given by chains of prime ideals, the arithmetical equivalent of the fact that on surfaces there are curves, and on curves there are points. This arithmetical formulation, which is proved identical with the parameter definition of elimination theory, can be transferred with a slight modification to arbitrary ring domains and, together with the transfer of certain parts of the other questions, gives there an exact insight into the structure of ideals, as I shall show elsewhere.

In H.-N. the question of \emph{zeroes} is, as usual, attacked directly for the given ideal, namely by going back to successive elimination, whereby the character of verification is not entirely avoided. In contrast stands the route always taken for polynomials in one variable: one represents the given polynomial as a product of powers of prime functions (primary functions) and constructs, for each individual prime function, the field of zeroes as a field isomorphic to the residue-class field. The degree of the prime function gives the degree of this field; the degree of the primary function, however, whose zeroes coincide with those of the prime function, gives the degree of the ring of residue classes modulo this primary function. If, for each prime function, one passes to the Galois field, one obtains the decomposition into linear factors; and with this, for the originally given polynomial, the question of zeroes, decomposition into linear factors, and multiplicity is settled. Quite correspondingly, here (§ 3) the system of residue classes modulo a prime ideal is extended, by forming quotients, to the residue-class field, and a field of zeroes is constructed that is isomorphic to it. This field of zeroes contains a subfield generated by adjoining indeterminates -- say $x_{i+1},\ldots,x_n$; the degree of the norm gives the degree with respect to this subfield; at the same time, by passage to the Galois field, the decomposition of the elementary-divisor form and hence of the norm into linear factors is obtained. For the corresponding primary ideal, however, the zeroes are the same; the decomposition is also given along with it, since the norm becomes a power of the elementary-divisor form of the prime ideal (§ 2); the degree of the norm gives the degree of the ring of residue classes modulo the primary ideal with respect to the subfield derived from the indeterminates.

Thus the question of zeroes of an arbitrary ideal is reduced to the connection between the decomposition theorems for norms and elementary divisors and those of general ideal theory (§ 5). The fact that the zeroes of the ideal are composed from those of its individual associated prime ideals corresponds to the decomposition of the norm into greatest primary factors, which become equal to powers of the elementary-divisor form of the individual associated prime ideals;\footnote{This parallelism between elimination theory and general ideal theory is not fulfilled in Kronecker's elimination theory; compare H.-N., note 4).} thereby the decomposition into linear factors is accomplished for the norm in general. Multiplicity, however, receives its interpretation in H.-N. by means of the fundamental ideals and their components; the degree of a greatest primary factor of the norm is made equal to the number of linearly independent residue classes -- after adjoining $x_{i+1},\ldots,x_n$ -- of the fundamental ideal of $(i-1)$-st stage modulo a component of the fundamental ideal of $i$-th stage. Now the fundamental ideal of $(i-1)$-st stage is proved identical with the isolated component of the decomposition that is uniquely defined by all prime ideals of dimension higher than the $(n-i)$-th; and a component of the fundamental ideal of $i$-th stage is proved identical with the isolated component of the decomposition defined by the prime ideal corresponding to the factor of the norm and by all prime ideals of higher dimension. The degree of the corresponding factor of the norm becomes -- after adjoining $x_{i+1},\ldots,x_n$ -- equal to the degree of that subring of the residue classes of this component that consists only of classes of the fundamental ideal of $(i-1)$-st stage.

The characterization of the prime ideals and proper primary ideals (§ 6) results as a direct consequence of the consideration of the residue-class fields. If the original coefficient domain is a perfect field, then prime functions correspond to the prime ideals as norm and elementary-divisor form, proper primary functions to the proper primary ideals, and conversely. For an imperfect field as coefficient domain only conditions can be given that are either necessary or sufficient; as examples show, the norm of a prime ideal can become properly primary, and the elementary-divisor form of a proper primary ideal can become a prime function. More precisely, in characteristic $p$ the norm of a prime ideal becomes the $p^g$-th power ($g\geqq0$) of its elementary-divisor form; whereas, for proper primary ideals whose elementary-divisor form is a prime function, the norm can become an arbitrary power, as examples again show. Finally, absolute prime ideals are considered (§ 7), that is, those that remain prime ideals when the coefficient domain is extended to an algebraically closed field. For absolute prime ideals with coefficients from a (finite) algebraic number field, the characterization leads to the transfer of a theorem first stated by A. Ostrowski for absolute prime functions: the property of being a prime ideal is preserved modulo every prime ideal of the number field, with at most finitely many exceptions.

Let the analogy of the last question with ideal theory in algebraic number fields also be pointed out. As the elementary-divisor form of such an ideal one has to regard the least integral rational number divisible by the ideal (note 10), which for a prime ideal therefore always becomes a prime number, whereas conversely an ideal becomes a prime ideal as soon as its norm is a prime number; the proofs also run quite in parallel. The examples mentioned above thus show that, over an imperfect field, the analogue of prime ideals of higher degree and ramification ideals occurs, whereas over a perfect field there are only prime ideals of first degree and no ramification ideals.

\subsection*{§ 1. Basic Concepts of Elimination Theory, General Ideal Theory, and Field Theory}

\textbf{1. The domain.} Let $\bar{\frakS}$ denote the integral domain of all polynomials in $y_1,\ldots,y_n$ with coefficients from an abstractly defined field $P$. Let the $y$ be subjected to a transformation with indeterminates $u_{\mu\nu}$ as coefficients:\footnote{The transformation is, with a view to § 3 and so on, somewhat more general than in H.-N.; all results there are thereby retained all the more.}
\srcnumdisplay{(1)}{y_1=u_{11}x_1+\cdots+u_{1n}x_n;\ \ldots;\ y_n=u_{n1}x_1+\cdots+u_{nn}x_n}
with inverse
\srcnumdisplay{(1')}{x_1=t_{11}y_1+\cdots+t_{n1}y_n;\ \ldots;\ x_n=t_{1n}y_1+\cdots+t_{nn}y_n,}
and the $u_{\mu\nu}$, or -- what is equivalent because of mutual rational expressibility -- the $t_{\mu\nu}$, are adjoined to $P$, so that the coefficient domain $P(u)=P(t)$ arises; let $\frakS$ denote the integral domain of all polynomials in $x_1,\ldots,x_n$ with coefficients from $P(u)$. The domains $\frakS$ and $\bar{\frakS}$ underlie elimination theory; ideals $\fraka,\frakb,\ldots$ in $\frakS$ and $\bar{\fraka},\bar{\frakb},\ldots$ in $\bar{\frakS}$ are considered. An ideal in an arbitrary ring is here defined in the usual way by the requirement that, together with $\alpha$ and $\beta$, it also contain $\alpha-\beta$, and, together with $\alpha$, also $\lambda\alpha$, where $\lambda$ is an arbitrary ring element; $\frako$ will always denote the unit ideal consisting of all elements.

\textbf{2. Transformed ideals.} An ideal $\frakm$ in $\frakS$ is called transformed if it arises from an ideal $\barfrakm$ in $\bar{\frakS}$ by means of (1) after adjoining the $u_{\mu\nu}$, or the $t_{\mu\nu}$. Thus $\frakm$ consists, when $\bar f(y)$ or $\bar g(y)$ runs through all polynomials in $\barfrakm$, of all linear combinations
\[
 f(x)=\sum U_i(u)\bar f_i(y)=\sum U_i(u)f_i(x)
\]
or also
\[
 g(x)=\sum T_i(t)\bar g_i(y)=\sum T_i(t)g_i(x);
\]
in particular all $f_i(x)=\bar f_i(y)$ are contained, or, what amounts to the same thing, all $g_i(x)=\bar g_i(y)$. Since $f(x)$ respectively $g(x)$ are fixed only up to quantities from $P(u)=P(t)$, the $U_i,T_i$ may, without loss of generality, be assumed to be power products of the $u$, respectively the $t$. The polynomials $f(x),g(x)$ again pass into polynomials of $\frakm$ if $U,T$ are replaced by arbitrary other power products, in particular by interchanging the columns of (1), respectively the rows of (1'); hence $\frakm$ contains, together with any polynomial, all those arising from it by interchanging the $x$, provided only that the corresponding interchanges are carried out in the coefficients depending on $u$, respectively $t$. All ideals occurring in what follows are to be assumed transformed unless the contrary is expressly stated. Non-transformed ideals pass, by composing (1) with $x_i=v_{i1}z_1+\cdots+v_{in}z_n$, into transformed ideals in the polynomial domain of the $z$ with coefficients from $P(u,v)$, whereas for transformed ideals this composition amounts merely to replacing the $u_{\mu\nu}$ by bilinear combinations of the $u,v$.

\textbf{3. Notation.} By $f^{(i)},g^{(i)},\ldots,a^{(i)},b^{(i)},\ldots$ we shall throughout understand polynomials in $\frakS$ that are free of $x_1,\ldots,x_{i-1}$.

\textbf{4. Fundamental ideals.} To every ideal $\frakm$ there are assigned $n$ fundamental ideals $\frakg_0,\frakg_1,\ldots,\frakg_{n-1}$ by the following stipulation: the fundamental ideal of $(i-1)$-st stage, $\frakg_{i-1}$, contains all and only those polynomials $G(x)$ for which there exists a polynomial $b^{(i)}\ne 0$ -- in general varying with $G(x)$ -- such that $b^{(i)}G(x)\equiv0(\frakm)$. From the existence of the ideal basis follows also the existence of at least one $B^{(i)}$, fixed for all $G(x)$, so that
{\let\veqno\leqno\begin{equation}
 B^{(i)}\frakg_{i-1}\equiv0(\frakm),\qquad B^{(i)}\ne0 . \tag{2}
\end{equation}}
The fundamental ideals of transformed ideals themselves become transformed ideals (H.-N., Theorem VI). The fundamental ideal of $i$-th stage $\frakg_i$ is also defined as the ideal into which $\frakm$ passes when $x_{i+1},\ldots,x_n$ are adjoined to $P(u)$, provided one restricts oneself -- which then means no loss of generality -- to polynomials integral in $x_{i+1},\ldots,x_n$; $\frakg_0$ is always equal to the unit ideal $\frako$.

\textbf{5. Module representation of ideals, elementary divisors and individual norms.} If every polynomial $f(x)$ is regarded as a linear form in the power products $\xi$ of $x_1,\ldots,x_{i-1}$, with polynomials $a^{(i)}$ as coefficients, then the ideal $\frakm$ passes into a module $\frakM_{i-1}$ of linear forms with respect to the domain of the $a^{(i)}$; that is, $\frakM_{i-1}$ contains, together with any two linear forms, also their difference, and, together with a linear form $l(\xi)$, also $a^{(i)}l(\xi)$ for arbitrary $a^{(i)}$. As fundamental module of such a module $\frakM$ one denotes the totality of linear forms $g(\xi)$ for which $b^{(i)}g(\xi)=0(\frakM)$, $b^{(i)}\ne0$; hence the fundamental ideal $\frakg_{i-1}$ passes, under the module representation, into the fundamental module $\frakG_{i-1}$ of $\frakM_{i-1}$. $\frakM_{i-1}$ and $\frakG_{i-1}$ depend on infinitely many indeterminates $\xi$, in such a way that in each individual linear form only finitely many of these indeterminates occur. $\frakM_{i-1}$ has the characteristic property that, after adjoining $x_{i+1},\ldots,x_n$ to $P(u)$, it has only finitely many elementary divisors different from unity; that is, after this adjunction $\frakG_{i-1}$ and $\frakM_{i-1}$ admit a basis representation -- with $\zeta$ denoting new indeterminates, connected with the $\xi$ by invertible linear transformations integral in $x_i$, in such a way that in $\zeta_i=r_i(\xi)$ only finitely many of these indeterminates occur at a time --
{\csname tagsleft@true\endcsname\begin{align}
\frakG_{i-1}&=(\zeta_0,\zeta_1,\ldots,\zeta_\sigma,\zeta_{\sigma+1},\ldots,\zeta_{\sigma+\nu},\ldots), \notag\\
\frakM_{i-1}&=(E^{(i)}\zeta_0,E_1^{(i)}\zeta_1,\ldots,E_\sigma^{(i)}\zeta_\sigma,
\zeta_{\sigma+1},\ldots,\zeta_{\sigma+\nu},\ldots), \tag*{\raisebox{9.5pt}[0pt][0pt]{(3)}}
\end{align}}
where each $E_\mu^{(i)}$ divides the preceding one (H.-N. (24) and (28)).

The highest elementary divisor $E^{(i)}$ is also defined as the greatest common divisor -- in the polynomial sense -- of all $B^{(i)}$ occurring in (2); it can be assumed integral and primitive in $x_{i+1},\ldots,x_n$, and is then regular in $x_i$, that is, $E^{(i)}$ contains the term $x_i^\lambda$ with non-zero coefficient, if it is of degree $\lambda$ in the $x$.

The product $R^{(i)}$ of all elementary divisors occurring in (3) is called the norm of $\frakG_{i-1}$ with respect to $\frakM_{i-1}$, or also the individual norm of the ideal $\frakm$; in symbols, taking account of the fact that $\frakm$ passes into $\frakg_i$ when $x_{i+1},\ldots,x_n$ are adjoined:
\begin{equation}
 R^{(i)}=E^{(i)}E_1^{(i)}\cdots E_\sigma^{(i)}
 =N(\frakG_{i-1}\mid\frakM_{i-1})=N(\frakg_{i-1}\mid\frakg_i).
\end{equation}
As (3) shows, after adjoining $x_{i+1},\ldots,x_n$, $R^{(i)}$ becomes equal to the determinant of the transition substitution from $\frakG_{i-1}$ to $\frakM_{i-1}$, and the degree of $R^{(i)}$ gives the number of residue classes, linearly independent over $P(u,x_{i+1},\ldots,x_n)$, of $\frakG_{i-1}$ modulo $\frakM_{i-1}$, hence of $\frakg_{i-1}$ modulo $\frakg_i$; $R^{(i)}$ can be assumed integral and primitive in $x_{i+1},\ldots,x_n$, and is then regular in $x_i$. $R^{(i)}$ can be computed as the greatest common divisor -- in the polynomial sense -- of all $\varrho$-rowed determinants of a module $\frakM^*_{i-1}$, always existing, of rank $\varrho$, having the property that the residue-class system $\frakG^*_{i-1}/\frakM^*_{i-1}$ is isomorphic to $\frakG_{i-1}/\frakM_{i-1}$, where $\frakG^*_{i-1}$ is understood as the fundamental module of $\frakM^*_{i-1}$ (H.-N., Theorem VII and § 5).

\textbf{6. Elementary-divisor form and norm (resultant form) of ideals.} The product $E_\frakm=E^{(1)}E^{(2)}\cdots E^{(n)}$ of all highest elementary divisors is called the elementary-divisor form, and the product $R_\frakm=R^{(1)}R^{(2)}\cdots R^{(n)}$ of all individual norms the norm (resultant form)\footnote{In H.-N. only the term resultant form is used; the term norm corresponds to the arithmetical properties.} of $\frakm$. The elementary-divisor form and the norm are divisible by the ideal; the norm is divisible by the elementary-divisor form, and a power of the latter by the norm. More generally one still has:
{\csname tagsleft@true\endcsname\begin{align}
 E^{(i)}\frakg_{i-1}&=0(\frakg_i),
& E^{(n)}\cdots E^{(i)}\frakg_{i-1}&=0(\frakm),\quad\text{and consequently}\notag\\
 R^{(i)}\frakg_{i-1}&=0(\frakg_i),
& R^{(n)}\cdots R^{(i)}\frakg_{i-1}&=0(\frakm),\quad\text{(H.-N., Theorem VIII),} \tag*{\raisebox{8.5pt}[0pt][0pt]{(4)}}
\end{align}}
and further: if $\frakm$ is divisible by $\frakn$, and the norms $R_\frakm$ and $R_\frakn$ agree, then \emph{the ideals $\frakm$ and $\frakn$ also agree} (H.-N., Theorem IX). Taking into account that, for a divisor of $\frakm$, agreement of $R^{(1)},R^{(2)},\ldots$ entails agreement of the fundamental ideals $\frakg_1,\frakg_2,\ldots$, and that always $\frakg_0=\frako$, one obtains correspondingly: if $\frakn$ is a \emph{proper divisor} of $\frakm$, then the first individual norm of $\frakn$ that differs from the corresponding one of $\frakm$ becomes a \emph{proper divisor}.\footnote{On the other hand, later factors of $R_\frakn$ can become multiples of the corresponding factors of $R_\frakm$; for instance always when $\frakm$ is a prime ideal and $\frakn$ is not the unit ideal.} \emph{If $\frakm$ contains a polynomial $G^{(i)}\ne0$}, then by definition $\frakg_0,\frakg_1,\ldots,\frakg_{i-1}$ are equal to the unit ideal; consequently, by the property of individual norms stated in 5, namely to determine the number of linearly independent residue classes, $R^{(1)},\ldots,R^{(i-1)}$ become equal to \emph{unity}; hence also $E^{(1)},\ldots,E^{(i-1)}$ become equal to \emph{unity}.

\textbf{7. Decomposition of the individual norms and elementary divisors; components of the fundamental ideals.} By a prime function one understands a polynomial irreducible with respect to $P(u)$; by a primary function, the power of a prime function. A proper primary function is one that is not at the same time a prime function, where therefore powers higher than the first are involved. A decomposition of a polynomial into greatest primary factors is one in which every factor is a primary function, but the product of any two factors is no longer primary.\footnote{The definition of prime function and primary function could also be formulated in exact analogy with 11 by replacing only the word ideal by polynomial. The greatest primary factors are the analogue of the greatest primary components of the decomposition in 12.} To the decomposition of the individual norm $R^{(i)}=N(\frakg_{i-1}\mid\frakg_i)$ into greatest primary factors $Q_\lambda^{(i)}$ there corresponds a representation of $\frakg_i$ as a least common multiple
\[
\frakg_i=[\frakr_1,\frakr_2,\ldots,\frakr_s],
\]
such that $Q_\lambda^{(i)}=N(\frakg_{i-1}\mid\frakr_\lambda)$; here $\frakr_\lambda$ possesses $\frakg_{i-1}$ as its fundamental ideal of $(i-1)$-st stage, and agrees with its fundamental ideal of $i$-th stage; the $\frakr_\lambda$ are called the components of the fundamental ideals. The highest elementary divisor of $\frakr_\lambda$ becomes equal to the greatest common divisor -- in the polynomial sense -- of $Q_\lambda^{(i)}$ and $E^{(i)}$, hence equal to a greatest primary factor of $E^{(i)}$. The $\frakr_\lambda$ are uniquely determined by their behavior, stated above, with respect to the fundamental ideals and by the requirement that either the individual norm or the highest elementary divisor should be a greatest primary factor of $R^{(i)}$ respectively $E^{(i)}$ (H.-N., Theorem X).

\textbf{8. Concept of dimension.} To the ideal $\frakm$ the highest dimension $n-i$ is assigned if $R^{(i)}$ is the first individual norm different from unity -- or, what by 6 is equivalent, if $\frakg_i$ is the first fundamental ideal different from the unit ideal; the dimension $-1$ is assigned to the unit ideal $\frako$. If there is only one individual norm $R^{(i)}$ different from unity, so that $\frakg_i=\frakg_{i+1}=\cdots=\frakm$, then the highest dimension is also called simply the dimension of $\frakm$. If $\frakm$ contains a polynomial $G^{(i)}\ne0$, then it is of highest dimension at most $n-i$, as the remark at the end of 6 shows. If $\frakm$ has highest dimension $n-i$ and $\frakn$ is a divisor of $\frakm$, then $\frakn$ also has highest dimension at most $n-i$.

\textbf{9. Norm (resultant form) and zeroes.} By zeroes of an ideal $\frakm$ one understands all systems of values of $x_1,\ldots,x_i$ belonging to a suitable algebraic extension field of $P(u;x_{i+1},\ldots,x_n)$ ($i=1,2,\ldots,n$) for which -- after adjoining $x_{i+1},\ldots,x_n$ to the coefficient domain -- all polynomials from $\frakm$ vanish. To each factor $R^{(i)}$ of the norm (resultant form), under this adjunction, there correspond as many zeroes of $\frakm$ as the degree of $R^{(i)}$ indicates; and $R^{(i)}$, written in bilinear combinations $w_{\mu\nu}$ of the $u_{\mu\nu}$ and further indeterminates $v$, admits the explicit decomposition:
\begin{equation}
 R^{(i)}=\prod\{z-(v_1\bar{x}_{1\nu}+\cdots+v_{i-1}\bar{x}_{i-1,\nu}+v_i\bar{x}_{i,\nu})\}^{\lambda_\nu}, \notag
\end{equation}
where $\bar{x}_{1\nu},\ldots,\bar{x}_{i,\nu}$ denotes a system of associated zeroes (H.-N., Theorem XIII; a new foundation will be given in § 3). Thus among the zeroes of an ideal of highest dimension $n-i$ there are zeroes depending on $(n-i)$ parameters, but none depending on more parameters; this proves the agreement of the concept of dimension given in 8 with the usual formulation.

\textbf{10. The ring domain of general ideal theory.} The underlying domain is to be a commutative ring in which the \emph{theorem of the finite chain} holds: every chain of ideals $\fraka_1,\fraka_2,\ldots,\fraka_k,\ldots$, where each $\fraka_i$ is a proper divisor -- divisor in the ideal sense -- of the immediately preceding one, breaks off after finitely many terms. This requirement is equivalent to the other, that every ideal possesses an ideal basis (Ideal Theory, Theorem I), and is therefore fulfilled in particular for the polynomial domain. In the following numbers 11, 12, 13, this general ring domain is assumed.

\textbf{11. Prime ideals and primary ideals.} An ideal $\frakp$ is called a prime ideal if from the divisibility of a product by $\frakp$ follows the divisibility of at least one factor; $\frakq$ is called a primary ideal -- or primary -- if from the divisibility of a product by $\frakq$ follows the divisibility of one factor or of a power of every factor. In symbols:
\[
 a\not\equiv0(\frakp),\quad b\not\equiv0(\frakp)\quad\hbox{implies}\quad ab\not\equiv0(\frakp);
\]
\[
 a\not\equiv0(\frakq),\quad b^x\not\equiv0(\frakq)\ \hbox{for every power }x\quad\hbox{implies}\quad ab\not\equiv0(\frakq).
\]
To every primary ideal $\frakq$ there is one and only one associated prime ideal $\frakp$, which is a divisor of $\frakq$ and a power of which is divisible by $\frakq$, $\frakq\equiv0(\frakp)$, $\frakp^\rho\equiv0(\frakq)$; here $\frakp$ consists of all ring elements of which a power is divisible by $\frakq$. As the definition shows, every prime ideal is at the same time a primary ideal; proper primary ideals mean those that are not at the same time prime ideals, so that $\rho>1$.

\textbf{12. The decomposition theorem.} A representation $\fraka=[\frakq_1,\ldots,\frakq_s]$ of an ideal as a least common multiple is called shortest if no $\frakq$ is contained in the least common multiple of the others -- in its complement; it is called a representation by greatest primary components if every $\frakq$ is primary, but the least common multiple of any two $\frakq$'s is not primary. Every ideal admits a shortest representation by \emph{finitely many greatest primary components}; for two different such representations, \emph{the number of components and the associated prime ideals, all of which are mutually distinct, agree}. The prime ideals thus uniquely determined are to be called the prime ideals, or associated prime ideals, of the ideal (Ideal Theory, Theorem IX).

\textbf{13. Isolated components of the decomposition.} An ideal $\frakr=[\frakq_{i_1},\ldots,\frakq_{i_\lambda}]$ is called an isolated component of the decomposition of $\fraka$ if the $\frakq_{i_j}$ occur in at least one shortest representation of $\fraka$ by greatest primary components and satisfy the condition that none of their associated prime ideals is contained in one of the other prime ideals of $\fraka$. The \emph{isolated components} of the decomposition are \emph{uniquely determined by their prime ideals}; in particular the isolated greatest primary components are uniquely determined (Ideal Theory, Theorem XIII).

\textbf{14. Characteristic, perfect and imperfect fields, extension of the first kind.} A field $\Omega$ is said to have characteristic zero if the prime field contained in $\Omega$ and derived from the unit is of the type of the field of rational numbers; it has characteristic $p$ if this prime field is of the type of the residue-class system modulo a prime number $p$. A field $\Omega$ is called perfect if every prime function with coefficients from $\Omega$ decomposes in a suitable extension field into distinct linear factors; otherwise it is called imperfect. Every field of characteristic zero is perfect; a field of characteristic $p$ is perfect if and only if, together with every element, it also contains its $p$-th root. A prime function in an imperfect field is an integral function of degree $r$ in $x^{p^f}$, and decomposes in a suitable extension field into $p^f$-th powers of $r$ distinct linear factors; $f$ is called the exponent. A prime function is said to be of the first kind if it decomposes in a suitable extension field into distinct linear factors, so that the exponent $f$ is zero; an extension is called of the first kind if every element is of the first kind, that is, is a zero of a prime function of the first kind. Every extension that arises by adjoining an element of the first kind is of the first kind; over perfect fields there are only extensions of the first kind (Steinitz, § 11 and § 13).

\textbf{15. Reduced degree and exponent of a finite extension.} If $\frK$ is a finite extension of an imperfect field $\Omega$, then $\frK$ contains a subfield $\frK_0$ of degree $r$ that is of the first kind with respect to $\Omega$ and consists of all and only the elements of the first kind in $\frK$. The $p^f$-th power of every element of $\frK$ belongs to $\frK_0$, and there are elements in $\frK$ that are exactly of degree $rp^f$. $r$ is called the reduced degree, $f$ the exponent of the finite extension (Steinitz, § 14, 1 and § 14, 5).

\subsection*{§ 2. Elementary-Divisor Form and Norm of the Prime Ideals and Primary Ideals. Divisors of the Prime Ideals}

From now on we shall throughout be concerned with the polynomial domain defined in § 1, 1. First it is to be shown that, also for prime ideals and primary ideals, one may restrict oneself to transformed ideals, according to

\textbf{Lemma I.} \emph{The transformed ideal $\frakp$ of a prime ideal $\bar{\frakp}$ in $\bar{\frakS}$ becomes a prime ideal; the transformed $\frakq$ of a primary ideal $\bar{\frakq}$ in $\bar{\frakS}$ becomes a primary ideal. In other words, the property of being a prime ideal or a primary ideal is preserved when the $u_{\mu\nu}$ are adjoined to $P$.}

Indeed suppose that $\fraka\not\equiv0(\frakp)$, $\frakb\not\equiv0(\frakp)$, where $\fraka$ and $\frakb$ need not be transformed ideals; say $a(x)\not\equiv0(\frakp)$ and $b(x)\not\equiv0(\frakp)$, with $a(x)$ and $b(x)$ polynomials from $\fraka$ and $\frakb$. By inversion of (1) according to (1') one obtains -- with $T$ understood, without loss of generality, as power products of the $t_{\mu\nu}$ --
\[
 a(x)=\sum T_j \bar a_j(y),\qquad b(x)=\sum T_j \bar b_j(y),
\]
and at least one $\bar a_j(y)$ and one $\bar b_j(y)$ must fail to be divisible by $\bar{\frakp}$. Let, for example, $\bar a_r(y)\not\equiv0(\bar{\frakp})$ and $\bar b_s(y)\not\equiv0(\bar{\frakp})$ be the respective highest terms under some lexicographic ordering of the $T$'s; then also $\bar a_r(y)\bar b_s(y)\not\equiv0(\bar{\frakp})$, and consequently $a(x)b(x)\not\equiv0(\frakp)$ and hence $\fraka\frakb\not\equiv0(\frakp)$, proving that $\frakp$ is prime. The proof evidently rests on the fact that the residue-class system modulo a prime ideal forms a ring without zero divisors, a property preserved by adjoining indeterminates.

Correspondingly suppose that $\fraka\not\equiv0(\frakq)$ and $\frakb^z\not\equiv0(\frakq)$ for every power $z$; then from the existence of the ideal basis it follows that there must exist an $a(x)$ from $\fraka$ and a $b(x)$ from $\frakb$ such that $a(x)\not\equiv0(\frakq)$ and $b(x)^z\not\equiv0(\frakq)$ for every power $z$; for under the contrary assumption, a power of $\frakb$ would be divisible by $\frakq$.

Putting $a(x)b(x)=c(x)$, let $\fraka^*,\frakb^*,\frakc^*$ denote the ideals respectively derived from the coefficients $\bar a_j(y),\bar b_j(y),\bar c_j(y)$ of $a(x),b(x),c(x)$. Then $\frakb^{*z}\not\equiv0(\bar{\frakq})$ for every power $z$, since otherwise $b(x)^z$ would be divisible by $\frakq$; because $\fraka^*\not\equiv0(\bar{\frakq})$, one must also have $\fraka^*\frakb^{*\lambda}\not\equiv0(\bar{\frakq})$ for every power $\lambda$. But by the Dedekind-Mertens theorem\footnote{Compare H.-N., p. 63 and note 13).} there exists an exponent $h$ such that $\fraka^*\frakb^{*h}=\frakb^{*h-1}\frakc^*$; hence $\frakc^*\not\equiv0(\bar{\frakq})$. This gives $c(x)\not\equiv0(\frakq)$ and consequently $\fraka\frakb\not\equiv0(\frakq)$, proving that $\frakq$ is primary.

Also in representing an ideal as a least common multiple one may restrict oneself to transformed ideals, according to

\textbf{Lemma II.} \emph{From a representation $\barfrakm=[\bar{\frakc}_1,\ldots,\bar{\frakc}_\alpha]$ in $\bar{\frakS}$ follows a representation $\frakm=[\frakc_1,\ldots,\frakc_\alpha]$ in $\frakS$, where $\frakc_i$ denotes the transformed ideal of $\bar{\frakc}_i$. In particular, therefore, in every shortest representation by greatest primary components (§ 1, 12.) these may be assumed transformed.}

Indeed $\frakm$ is divisible by $[\frakc_1,\ldots,\frakc_\alpha]$, since every polynomial $f(x)$ from $\frakm$ -- after multiplication, if necessary, by a suitable power of $u_{\mu\nu}$ -- is of the form $\sum U_i\bar f_i(y)$, where each $\bar f_i(y)$, as a polynomial from $\barfrakm$, is by assumption divisible by all the $\bar{\frakc}_i$. Conversely, if $f(x)$ is divisible by every $\frakc_i$, then it has the above form and is therefore divisible by $\frakm$. Thus if one starts from a decomposition into greatest primary components in $\bar{\frakS}$, one obtains a decomposition $\frakm=[\frakq_1,\ldots,\frakq_\alpha]$ in $\frakS$; here the $\frakq$, as transformed ideals of primary ideals, are primary by Lemma I, and they are greatest primary components, since distinct associated prime ideals $\bar{\frakp}$ in $\bar{\frakS}$ correspond to distinct $\frakp$ in $\frakS$.

We now need a first characterization of prime ideals and primary ideals by elementary-divisor form and norm, still without complete separation of prime and primary. In exact analogy with ideal theory in algebraic number fields,\footnote{As the highest elementary divisor of ideals in algebraic number fields one has to regard the least integral rational number divisible by the ideal, hence in particular the prime number divisible by a prime ideal or the prime-power divisible by a primary ideal (a power of a prime ideal). In fact every ideal $\fraka^*$ is also a module $A^*$ of linear forms in the elements $\omega_1,\ldots,\omega_n$ of a field basis, with $(\omega_1,\ldots,\omega_n)$ the fundamental module of $A^*$; the least integral rational number divisible by $\fraka^*$ therefore becomes the highest elementary divisor of this module $A^*$, whereas the norm of $\fraka^*$ becomes the product of all elementary divisors of $A^*$.}

\textbf{Theorem I.} \emph{The elementary-divisor form of a prime ideal becomes equal to a prime function, and the norm to a power of this prime function. For primary ideals, elementary-divisor form and norm become primary functions, namely powers of the same prime function, which itself is the elementary-divisor form of the associated prime ideal. For prime ideals and primary ideals, highest dimension coincides with dimension simpliciter; this dimension agrees for a primary ideal and its associated prime ideal.}

Theorem I is evidently fulfilled for the unit ideal, whose elementary-divisor form and norm become unity. Now let the prime ideal $\frakp$ have highest dimension $(n-i)$; by (4), § 1, 6, one obtains
\[
 E^{(n)}\cdots E^{(i+1)}\frakg_i\equiv0(\frakp),\quad\hbox{but}\quad
 E^{(n)}\cdots E^{(i+1)}\not\equiv0(\frakp),
\]
since otherwise $\frakp$ would, by § 1, 8, have highest dimension at most $n-i-1$. Hence $\frakg_i\equiv0(\frakp)$, so that $\frakp$ is identical with its fundamental ideal of $i$-th stage and consequently also with $\frakg_{i+1},\frakg_{i+2},\ldots$. Thus $R^{(i+1)}=N(\frakg_i\mid\frakg_{i+1})$, $R^{(i+2)},\ldots$ become equal to unity, and hence also $E^{(i+1)},E^{(i+2)},\ldots$, as divisors of $R^{(i+1)},R^{(i+2)},\ldots$; therefore the elementary-divisor form becomes equal to $E^{(i)}$, which must be equal to a prime function $P^{(i)}$. For by § 1, 5 and taking into account that $\frakg_{i-1}$ becomes the unit ideal, $E^{(i)}$ is defined as the greatest common divisor -- in the polynomial sense -- of all $B^{(i)}$ divisible by $\frakp$; from $E^{(i)}=C_1^{(i)}C_2^{(i)}\equiv0(\frakp)$ it would follow that $E^{(i)}$ must be contained in one of its factors $C_1^{(i)}$ or $C_2^{(i)}$. Since, however, a power of $E^{(i)}$ is divisible by $R^{(i)}$, $R^{(i)}$ becomes a power -- possibly the first power -- of this prime function $P^{(i)}$; at the same time $R^{(i)}$ becomes the norm of $\frakp$. Since only one individual norm different from unity occurs, the highest dimension of $\frakp$ finally agrees with its dimension simpliciter.

Correspondingly, let the primary ideal $\frakq$ have highest dimension $n-i$; then, as above,
\[
 E^{(n)}\cdots E^{(i+1)}\frakg_i\equiv0(\frakq),\quad
 (E^{(n)}\cdots E^{(i+1)})^z\not\equiv0(\frakq)
\]
for every power $z$; and consequently, as above, $\frakq=\frakg_i$, its elementary-divisor form is $E^{(i)}$, its norm is $R^{(i)}$, and its highest dimension is its dimension simpliciter. This dimension agrees with that of the associated prime ideal; for from $\frakp^\rho\equiv0(\frakq)$, $\frakq\equiv0(\frakp)$ it follows that the dimension $(n-i)$ of $\frakp$ is at most equal to that of $\frakq$, and that of $\frakq$ at most equal to that of $\frakp^\rho$; from $(P^{(i)})^\rho\equiv0(\frakp^\rho)$, where $P^{(i)}$ denotes the elementary-divisor form of $\frakp$, the latter highest dimension is found to be at most $n-i$, and thus $n-i$ is the dimension of $\frakp$ and $\frakq$. From $(P^{(i)})^\rho\equiv0(\frakq)$ it follows further that $E^{(i)}$ -- as the greatest common divisor of all $B^{(i)}$ from $\frakq$ -- becomes a power of $P^{(i)}$, which may also be the first power, and that the same holds for $R^{(i)}$. This proves Theorem I in all its parts.

A characterization of prime ideals by means of the concept of dimension is given by

\textbf{Theorem II.} \emph{A prime ideal has no proper divisor of the same highest dimension; conversely, every ideal with this property is prime.}

The proof rests on

\textbf{Lemma III.} \emph{If a prime ideal $\frakp$ of polynomials with coefficients from a field $\Omega$ has only finitely many residue classes linearly independent over $\Omega$, then $\frakp$ has no proper divisor different from the unit ideal; in other words, the residue-class system modulo $\frakp$ forms a field.}

The hypothesis says that there is a finite number $k$ such that among any $(k+1)$ residue classes there is a linear dependence with coefficients from $\Omega$ -- more precisely, with coefficients from the residue-class field $(\Omega)$ represented by all elements of $\Omega$ -- but that at least one system of $k$ residue classes linearly independent over $(\Omega)$ exists. It follows immediately that all residue classes can be expressed linearly by any chosen system of $k$ linearly independent ones. Now let $\fraka$ be a proper divisor of $\frakp$, and let $a(z)\not\equiv0(\frakp)$ be a polynomial from $\fraka$; from $c_1f_1(z)+\cdots+c_kf_k(z)\not\equiv0(\frakp)$ it follows that
\[
 c_1a(z)f_1(z)+\cdots+c_ka(z)f_k(z)\not\equiv0(\frakp);
\]
that means that, together with the residue classes of $f_1,\ldots,f_k$, the residue classes of $af_1,\ldots,af_k$ also form a system of $k$ linearly independent ones, through which in particular the unit class can be linearly represented. Hence $\fraka$ is the unit ideal; expressed differently, the residue-class system forms a field, since for $a(z)\not\equiv0(\frakp)$ there is always an $f(z)$ such that $1\equiv a(z)f(z)(\frakp)$.

For the proof of the first part of Theorem II it remains only to observe that every prime ideal of dimension $(n-i)$ -- more generally every ideal of highest dimension $(n-i)$ -- has only finitely many residue classes linearly independent over $P(u;x_{i+1},\ldots,x_n)$, namely, by § 1, 5, as many as the degree of $R^{(i)}$ indicates, since here $\frakg_{i-1}$ becomes the unit ideal. Every proper divisor $\fraka$ of $\frakp$ therefore becomes the unit ideal when $x_{i+1},\ldots,x_n$ are adjoined to $P(u)$; expressed differently, the fundamental ideal of $i$-th stage of $\fraka$ becomes the unit ideal, which means that the highest dimension of $\fraka$ is at most $n-i-1$. In order that the assertion also hold for a non-transformed ideal $\fraka$ as divisor, one has only to pass, according to § 1, 2, to a new transformed domain.

Conversely suppose now that $\frakp$ has no proper divisor of the same highest dimension $n-i$, and let $\fraka\not\equiv0(\frakp)$ and $\frakb\not\equiv0(\frakp)$, with $\fraka$ and $\frakb$ again assumed to be transformed ideals. Then, by hypothesis, since $(\fraka,\frakp)$ and $(\frakb,\frakp)$, as greatest common divisors -- in the ideal sense -- become proper divisors of $\frakp$, one has
\[
 G^{(i+1)}\equiv0(\fraka,\frakp),\qquad
 H^{(i+1)}\equiv0(\frakb,\frakp),\qquad
 G^{(i+1)}\not\equiv0,\qquad
 H^{(i+1)}\not\equiv0.
\]
This gives
\[
 G^{(i+1)}H^{(i+1)}\equiv0(\fraka\frakb,\frakp),\qquad
 G^{(i+1)}H^{(i+1)}\not\equiv0;
\]
and consequently necessarily $\fraka\frakb\not\equiv0(\frakp)$, since otherwise $\frakp$ would have smaller highest dimension than $n-i$. Since $\fraka,\frakb$ are transformed ideals, $\bar{\frakp}$, and by Lemma I also $\frakp$, are prime ideals. This proves Theorem II.

\subsection*{§ 3. Zeroes of the Prime Ideals and Primary Ideals}

The passage to a direct foundation of the theory of zeroes (§ 1, 9.), at first for prime ideals, is formed by the following lemma, which extends Lemma III and is valid for prime ideals in arbitrary rings:

\textbf{Lemma IV.} \emph{The residue-class system modulo every prime ideal different from the unit ideal can be extended by quotient formation to a field, the residue-class field of the prime ideal.}

The residue-class system modulo an arbitrary ideal forms a ring, since sum, difference and product again belong to the system, and the associative, commutative and distributive laws are preserved on passage to residue classes. If the ideal is specifically prime, this ring has no zero divisors; for the definition of prime ideals (§ 1, 11.) says that the product of non-vanishing residue classes likewise does not vanish. If the prime ideal is different from the unit ideal, then this ring contains at least one element different from the zero element, and can therefore be extended to a field by quotient formation -- adjunction of pairs of elements;\footnote{Steinitz, § 3. The assumption of \emph{one} non-zero element in the original integral domain suffices, since then the existence of the unit is secured by quotient formation, and hence the original domain becomes a subdomain of the field; Steinitz assumes the existence of the unit in the integral domain.} thus the residue-class field is defined.

We first fix the notation that will be kept throughout what follows:

\textbf{Notation:} Residue classes shall generally be denoted by putting any representative of the class in parentheses, $(x),\ldots,(a(x)),\ldots$; similarly, fields of residue classes shall be denoted by parentheses. In particular, $(P)$, $(P(u))$, $(P(u,x_{i+1},\ldots,x_n))$ denote the residue-class fields consisting of all and only those residue classes that can be represented by elements from $P$, $P(u)$, $P(u,x_{i+1},\ldots,x_n)$.

We also recall the following: A field $\mR$ is said to have \emph{algebraic rank (transcendence degree)} $k$ with respect to a base field $\Omega$ if $\mR$ contains at least one system of $k$ quantities algebraically independent over $\Omega$ -- transcendental quantities -- but every $k+1$ quantities from $\mR$ are algebraically dependent over $\Omega$. If $\mR$ can be represented as an algebraic extension field of a subfield $\Omega(z_1,\ldots,z_s)$, where the $z$ are algebraically independent -- transcendental -- over $\Omega$, then necessarily $s=k$;\footnote{For a proof of this familiar fact valid in arbitrary characteristic, compare Steinitz, § 22, 9. -- On adjoining the $u$, the algebraic rank cannot increase, since only rational combinations of the original elements with coefficients from $\Omega(u)$ are involved; nor can it decrease, since every relation between elements of $\mR$ with coefficients from $\Omega(u)$ decomposes into finitely many relations with coefficients from $\Omega$.} likewise the algebraic rank of $\mR(u)$ over $\Omega(u)$ is equal to $k$ when $u$ denotes indeterminates not contained in $\mR$.

The construction of the field of zeroes now proceeds in complete analogy with the usual route for prime functions of one indeterminate,\footnote{Steinitz, § 6.} according to

\textbf{Theorem III.} \emph{To the residue-class field $(\mR)$ of a prime ideal $\frakp$ of dimension $n-i$ there can be assigned, isomorphically, a field of zeroes $R$ of $\frakp$, which has algebraic rank $n-i$ -- depends on $n-i$ parameters -- and in particular can be regarded as a finite extension field of $P(u,x_{i+1},\ldots,x_n)$. The isomorphism between $(\mR)$ and $R$ includes that between $(P)$ and $P$, $(P(u))$ and $P(u)$, and, when $x_{i+1},\ldots,x_n$ are taken as parameters, also that between $(P(u,x_{i+1},\ldots,x_n))$ and $P(u,x_{i+1},\ldots,x_n)$. Conversely, every field of zeroes of algebraic rank $n-i$ -- depending on $n-i$ parameters -- is isomorphic to the residue-class field $(\mR)$.}

Theorem III yields as a corollary the familiar theorem: \emph{If an ideal $\fraka$ vanishes in all zeroes of a prime ideal $\frakp$, then $\fraka$ is divisible by $\frakp$.}

For the proof, first note that the residue-class field $(\mR)$ has algebraic rank $n-i$ with respect to $(P(u))$. Indeed, the classes $(x_{i+1}),\ldots,(x_n)$ are algebraically independent with respect to $(P(u))$, since $\frakp$, being of dimension $n-i$, contains no polynomial free of $x_1,\ldots,x_i$; every further class, however, is dependent on these. For from the membership of the elementary-divisor form $P^{(i)}$ in $\frakp$ it follows, by § 1, 2., also for $\lambda=1,2,\ldots,i$ that polynomials belong to $\frakp$ which depend respectively on $x_\lambda,x_{i+1},\ldots,x_n$. Thus $(\mR)$ becomes a finite extension field of $(P(u,x_{i+1},\ldots,x_n))$; by § 1, 5. -- taking into account that $\frakg_{i-1}$ is equal to the unit ideal and $\frakg_i$, by Theorem I, to $\frakp$ -- the degree over this subfield is equal to the degree of the norm.

To pass to an isomorphic field of zeroes, one further observes that $P(u,x_{i+1},\ldots,x_n)$ and $(P(u,x_{i+1},\ldots,x_n))$ are related isomorphically by assigning to each element of $P(u,x_{i+1},\ldots,x_n)$ the residue class represented by that element. First, under this assignment, the integral domains consisting of polynomials in $u,x_{i+1},\ldots,x_n$ and respectively in $(u),(x_{i+1}),\ldots,(x_n)$ correspond isomorphically, since $\frakp$ contains no polynomial free of $x_1,\ldots,x_i$, and therefore distinct polynomials from $P(u,x_{i+1},\ldots,x_n)$ correspond to distinct classes; from isomorphic integral domains, however, isomorphic fields arise by quotient formation. The isomorphic relation between $P(u,x_{i+1},\ldots,x_n)$ and $(P(u,x_{i+1},\ldots,x_n))$ includes that between $P$ and $(P)$, $P(u)$ and $(P(u))$. To extend this isomorphism to one between $(\mR)$ and a field of zeroes $R$, let certain new elements $\bar{x}_1,\ldots,\bar{x}_i$ be assigned isomorphically to the classes $(x_1),\ldots,(x_i)$; that is, to each polynomial $(a(x))$ there corresponds a polynomial $a(\bar{x}_1,\ldots,\bar{x}_i,x_{i+1},\ldots,x_n)$, and sums and products correspond to sums and products. By quotient formation -- where one may restrict oneself to denominators from the subfield $(P(u,x_{i+1},\ldots,x_n))$ respectively $P(u,x_{i+1},\ldots,x_n)$ -- to the residue-class field $(\mR)$ there then corresponds a field $R$, a finite extension of $P(u,x_{i+1},\ldots,x_n)$ of the degree of the norm; it can be called the field of zeroes. For the assignment says that $f(\bar{x}_1,\ldots,\bar{x}_i,x_{i+1},\ldots,x_n)$ vanishes as soon as $f(x)\equiv0(\frakp)$. In place of $(P(u,x_{i+1},\ldots,x_n))$, any other subfield of $(\mR)$ arising by adjoining $n-i$ algebraically independent quantities could occur throughout the argument.

Conversely it is easy to show that every field of zeroes $R$ of algebraic rank $n-i$ is isomorphic to the residue-class field. Since $R$ can be derived rationally, with coefficients from $P(u)$, from the quantities $\bar{x}_1,\ldots,\bar{x}_n$, among these quantities there must be $n-i$ algebraically independent, hence transcendental, elements with respect to $P(u)$. In particular this must hold for $\bar{x}_{i+1},\ldots,\bar{x}_n$, since, because $P^{(i)}\equiv0(\frakp)$, it follows as above that $R$ is an algebraic extension field of $P(u,\bar{x}_{i+1},\ldots,\bar{x}_n)$. Since by hypothesis $f(\bar{x})$ vanishes for every polynomial $f(x)$ in $\frakp$, to every class of the polynomial domain of the $(x)$ contained in $(\mR)$ there corresponds one and only one element of $R$ which is a polynomial in the $\bar{x}$; and sums and products correspond to sums and products. But different classes correspond to different elements of $R$; otherwise a non-zero class, and hence after suitable multiplication also a class represented by a polynomial from $(P(u,x_{i+1},\ldots,x_n))$, would correspond to the element zero in $R$, and contrary to the hypothesis the quantities $\bar{x}_{i+1},\ldots,\bar{x}_n$ would satisfy an algebraic dependence. This assignment exhausts all polynomials in the $\bar{x}$ contained in $R$; by quotient formation -- where one may again restrict oneself to denominators from $(P(u,x_{i+1},\ldots,x_n))$ respectively $P(u,\bar{x}_{i+1},\ldots,\bar{x}_n)$ -- isomorphic fields arise from these isomorphic integral domains.

Now let the ideal $\fraka$ vanish in all zeroes of $\frakp$, in particular also in all zeroes depending on $n-i$ parameters. If $a(x)$ is any polynomial from $\fraka$, then $a(\bar{x})$ vanishes; because of the isomorphism between $R$ and $(\mR)$, the class $(a(x))$ is thus the zero class, and $a(x)$, hence $\fraka$, is divisible by $\frakp$. This proves Theorem III and its corollary.

To pass from Theorem III to the decomposition given in § 1, 9., first for prime ideals, and in order to be able to make more precise statements about the exponents, a further investigation of the elementary-divisor form is needed, according to

\textbf{Lemma V.} \emph{If in characteristic $p$ the elementary-divisor form $E^{(i)}$ of a prime ideal or primary ideal -- more generally, of an ideal $\fraka$ for which highest dimension and dimension simpliciter coincide -- has exponent $f$ with respect to $x_i$, hence is an integral function of $x_i^{p^f}$, then it also has exponent $f$ with respect to $x_{i+1},\ldots,x_n$ and with respect to the indeterminates $u_{\mu r}$ respectively $t_{\mu r}$ occurring in $E^{(i)}$. In particular, if the coefficient domain is a perfect field, the elementary-divisor form $P^{(i)}$ of a prime ideal always has exponent zero with respect to $x_i$, hence is a prime function of the first kind.}

For this it is to be observed that $P(u,x_{i+1},\ldots,x_n)$ becomes imperfect as soon as $P$ is a perfect field of characteristic $p$, so that it does not follow directly from § 1, 14. that $P^{(i)}$ is of the first kind. Let $E^{(i)}$ have exponent $f$ with respect to $x_i$, but let it contain another indeterminate $x_{i+\lambda}$ with exponent $f'$. Then -- by interchanging the $u_{\mu r}$ respectively $t_{\mu r}$ -- there also exists, by § 1, 2., a non-zero polynomial $G^{(i)}$, divisible by the ideal, that is an integral function of $x_i^{p^{f'}}$, and which, by definition of the elementary-divisor form, is divisible by $E^{(i)}$. But since $G^{(i)}$ has the same degree as $E^{(i)}$, it must agree with $E^{(i)}$ up to a factor from $P(u)$; hence necessarily $f'=f$. Thus $E^{(i)}$ -- which may be assumed integral and primitive in the $t_{\mu r}$ -- is of the form
\[
  E^{(i)}=\sum c_\lambda(t)\,x_i^{\lambda_i p^f}\cdots x_n^{\lambda_n p^f}
       =\sum c_i(t)\,g_i(y^{p^f},t^{p^f})
       =\sum T_\mu(t)\,\bar h_\mu(y^{p^f}),
\]
where the $T_\mu$ are power products of the $t$, and the $\bar h$ are divisible by $\bar{\fraka}$. Therefore one obtains a polynomial again belonging to $\fraka$ if in the $T_\mu$ every exponent of the individual $t$ is replaced by the greatest multiple of $p^f$ contained in it; as the displayed representation shows, this amounts to replacing, in $c_\lambda(t)$, every exponent of the $t$ by the greatest multiple of $p^f$ contained in it. By this substitution $E^{(i)}$ passes into a polynomial in $t,x_i,\ldots,x_n$ which, by definition, is divisible by $E^{(i)}$, and, because of the assumed primitivity, also with respect to the $t$; hence it must be identical with $E^{(i)}$, since the degree has increased in no argument. In particular, if $P$ is perfect, then $E^{(i)}$, as an integral function of $x^{p^f},t^{p^f}$, is a $p^f$-th power; hence, if a prime ideal is involved, necessarily $f=0$, and the elementary-divisor form $P^{(i)}$ becomes a prime function of the first kind.

The decomposition is now given by

\textbf{Theorem IV.} \emph{The elementary-divisor form $P^{(i)}$ of a prime ideal admits, in the Galois field derived from the field of zeroes, the explicit decomposition:}
\[
 P^{(i)}=\prod (x_i-t_{1i}\bar y_{1\nu}-\cdots-t_{ni}\bar y_{n\nu})^\delta
 =\prod (t_{1i}(y_1-\bar y_{1\nu})+\cdots+t_{ni}(y_n-\bar y_{n\nu}))^\delta,
\]
\emph{where $\bar y_{1\nu},\ldots,\bar y_{n\nu}$ denotes in each case an associated system of zeroes of the original indeterminates $y$, independent of $t_{1i},\ldots,t_{ni}$; distinct $\nu$ correspond to distinct linear factors; and the exponent $\delta$ has the value one when $P$ is perfect and the value $p^f$ ($f\ge0$) when it is imperfect. Thereby, after adjoining new indeterminates $v$, the decomposition of the elementary-divisor form written in bilinear combinations of the $u$ and $v$ in place of the $u_{\mu\nu}$ is also determined:}
\[
 P^{(i)}(u,v)=\prod (z-v_1\bar x_{1\nu}-\cdots-v_i\bar x_{i\nu})^\delta,
\]
\emph{where $\delta$ has the same meaning as above. Likewise the decomposition of the norm of a prime ideal, or of a primary ideal with $\frakp$ as associated prime ideal -- as a power of $P^{(i)}$ -- is also given.}

As a corollary of Theorem IV one also obtains: \emph{If two prime ideals agree in their elementary-divisor form $P^{(i)}$, then they are identical.}

The proof of Theorem IV will amount to proving that the elementary-divisor form $P^{(i)}$ is identical with the fundamental equation of the extension field $(P(u,x_i,x_{i+1},\ldots,x_n))$ over $(P(u,x_{i+1},\ldots,x_n))$. First, to gain insight into the dependence on the indeterminates $u_{\mu r}$ respectively $t_{\mu r}$, pass to the residue-class field $(\bar{\mR})$ of $\bar{\frakp}$, which, by what was said in the definition of algebraic rank, must have algebraic rank $n-i$ with respect to $(\bar P)$, where $(\bar P)$ denotes the subfield represented by elements of $P$. For $(\bar{\mR})$ arises by adjoining the indeterminates $u$ respectively $t$ to a subfield isomorphic to $(\mR)$, obtained by quotient formation from all and only those classes representable by polynomials from $\bar{\frakS}$, with $(\bar P)$ and $(P)$ corresponding to one another. The algebraic rank of $(\bar{\mR})$ with respect to $(P(u))$ is therefore equal to that of this subfield with respect to $(\bar P)$, and hence, by the isomorphism, to that of $(\mR)$ with respect to $(P)$. Since $(\bar{\mR})$ can be derived rationally, with coefficients from $(\bar P)$, from the classes $(y_1),\ldots,(y_n)$, there must be among these classes $n-i$ algebraically independent ones, and $(\bar{\mR})$ arises as a finite algebraic extension field of the subfield generated by adjoining these $n-i$ classes to $(\bar P)$.\footnote{This remark shows that Theorem III with its corollary could have been stated and proved directly for the field of zeroes $\bar R$ of $\bar{\frakp}$. Only by passing to the transformed ideal is it achieved that in the greatest primary components of equal dimension of an arbitrary ideal the same parameters occur among the zeroes; and only thereby does the connection with norm and elementary-divisor form of an arbitrary ideal result.}

If one adjoins only the indeterminates $t_{\sigma\tau}$ occurring in $x_{i+1},\ldots,x_n$, then again a residue-class field arises which contains the classes $(y_1),\ldots,(y_n)$ and $(x_{i+1}),\ldots,(x_n)$ and becomes a finite algebraic extension field of $(P(t_{\sigma\tau},x_{i+1},\ldots,x_n))$; these are fields isomorphic to subfields of $(\bar{\mR})$, not identical with them, since for different indeterminates the residue classes are represented by the same elements but are not identical, because they do not contain the same elements. Thus in the dependence of the classes $(y)$ on $(P(t_{\sigma\tau},x_{i+1},\ldots,x_n))$ only the indeterminates $t_{\sigma\tau}$ enter. Now form, by adjoining $t_{1i},\ldots,t_{ni}$, the class $(x_i)=(t_{1i}y_1+\cdots+t_{ni}y_n)$; let $(x_{i\nu})=(t_{1i}y_{1\nu}+\cdots+t_{ni}y_{n\nu})$ be the $r$ distinct conjugate values lying in the Galois field with respect to $(P(t_{\sigma\tau},t_{1i},\ldots,t_{ni},x_{i+1},\ldots,x_n))$. Then -- according as one is dealing with extensions of the first kind or not -- the product over the factors $t-(x_{i\nu})$, or over their $p^f$-th powers, is a prime function $(G)$ with respect to $(P(t,x_{i+1},\ldots,x_n))$; for, since there are elements of degree $r\cdot p^f$ (§ 1, 15.), $(x_i)$ must necessarily be one such element, since every element arises by specializing the $t_{1i},\ldots,t_{ni}$; but $(x_i)$ is a zero of $(G)$. Passing to the isomorphic field of zeroes of $\frakp$ and to the Galois field derived from it, $G$ therefore admits a decomposition into linear factors $t-t_{1i}\bar y_{1\nu}-\cdots-t_{ni}\bar y_{n\nu}$ respectively their $p^f$-th powers. Now, since $(G)$ vanishes for $t=(x_i)$, $G(x_i)$ is divisible by $\frakp$ and hence by $P^{(i)}$; and, as a prime function with respect to $P(u,x_{i+1},\ldots,x_n)$ and as primitive in $x_{i+1},\ldots,x_n$, $G$ must become identical with $P^{(i)}$. Thus the first decomposition of Theorem IV is proved.

To pass to the second decomposition, note that the property of being, or not being, an extension of the first kind is preserved when indeterminates are adjoined. Therefore, after adjoining $v_1,\ldots,v_i$ and all the $t_{\mu\nu}$, form the class $(z)=(v_1x_1+\cdots+v_ix_i)$ and the conjugate classes $(z_\nu)=(v_1x_{1\nu}+\cdots+v_ix_{i\nu})$. Then the product over all $t-(z_\nu)$, or over their $p^f$-th powers, is again a prime function $(H)$ that vanishes for $t=(z)$. Hence $H$ is representable as a product of linear factors, as above, and is identical with the elementary-divisor form of the prime ideal derived from $\frakp$ by composing the transformation (1) with $z=v_1x_1+\cdots+v_ix_i$; this elementary-divisor form, however -- because of the mutual specialization -- is obtained from $P^{(i)}$ by replacing the $u_{\mu\nu}$ by certain bilinear combinations of the $u$ and $v$.\footnote{Compare H.-N., § 7.} With these decompositions, the decomposition of the norm as a power of $P^{(i)}$ is also given,\footnote{For extensions of the first kind, hence in particular over perfect fields, the first power is involved; over imperfect fields the $p^g$-th power is involved, as will be shown in § 6, Theorems XIII and XIV.} and likewise that of the norm and elementary-divisor form of a primary ideal with $\frakp$ as associated prime ideal.

Finally, if two prime ideals agree in their elementary-divisor form $P^{(i)}$, then, as a consequence of the above decomposition, they agree in their zeroes; hence they are mutually divisible by one another and are therefore identical.

\subsection*{§ 4. Arithmetical Formulation of the Concept of Dimension.}

A first formulation of the concept of dimension that is independent of the introduction of transformed ideals is given by

\textbf{Theorem V.} \emph{The dimension of a prime ideal is given by the algebraic rank of its residue-class field. The highest dimension of an ideal is equal to the highest of the dimension numbers of its associated prime ideals.}

The first part of Theorem V was proved in the first remark to Theorem IV, where it was shown that the algebraic rank of the residue-class field $(\bar{\mR})$ of $\bar{\frakp}$ is equal to that of the residue-class field $(\mR)$ of $\frakp$, hence equal to the dimension of $\frakp$, in agreement with § 1, 8. For the proof of the second part, let $\frakm=[\frakq_1,\ldots,\frakq_a]$ be a shortest representation by greatest primary components. Then, first of all, no $\frakq$ as a divisor of $\frakm$ can possess higher dimension than $\frakm$. But the product of all $\frakq$ is divisible by $\frakm$, and therefore so is the product of all elementary-divisor forms of the individual $\frakq$'s. If, then, $n-i$ is the highest of the dimension numbers of the $\frakq$, then $\frakm$ contains a polynomial $G^{(i)}\ne0$ and can therefore not have dimension higher than $n-i$. The dimension numbers of the $\frakq$ agree, by Theorem I, with those of their associated prime ideals. Thus if, without introducing transformed ideals, one defines the dimension of a prime ideal by the algebraic rank of its residue-class field, then the second part of Theorem V gives the definition of the highest dimension of an arbitrary ideal.

In order, on the other hand, to grasp the concept of dimension by chains of prime ideals -- again independently of the introduction of transformed ideals -- one has, supplementing Theorem II, to show:

\textbf{Theorem VI.} \emph{Every prime ideal different from the unit ideal possesses -- after possible adjunction of indeterminates -- at least one prime ideal as divisor whose dimension has decreased by exactly one unit.}

Indeed, if $\frakp$ has dimension $(n-i)>0$ and $v$ denotes a new indeterminate, then
\[
 \frakc=(\frakp,x_i-v)
\]
is an ideal with the required property. This follows from an isomorphic correspondence. One has $\frakc=(\frakp^*,x_i-v)$, where $\frakp^*$ arises from $\frakp$ by replacing $x_i$ by the indeterminate $v$ and adjoining $v$ to the coefficient domain; the residue classes are likewise obtained by replacing the class $(x_i)$ by $(v)$. But $\frakp$ passes into itself when $x_i$ is adjoined to $P(u)$; for from
\[
 h(x_i)f(x)\equiv0(\frakp),\qquad h(x_i)\ne0
\]
one necessarily obtains $f(x)\equiv0(\frakp)$. Otherwise $\frakp$ would contain a polynomial depending only on $x_i$, and therefore also one depending only on $x_n$, contrary to the supposition it would have dimension at most zero. Under the assignment $v\sim x_i$, therefore, the zero class modulo $\frakc$ necessarily corresponds to the zero class modulo $\frakp$, and of course conversely as well. Hence there is an isomorphism between the residue-class system modulo $\frakp$ and that modulo $\frakc$; the latter has no zero divisors, so $\frakc$ is a prime ideal. Under this isomorphism the algebraic dependence between $(x_i)$ and $(P(u,x_{i+1},\ldots,x_n))$ corresponds to a relation between $(x_{i+1}),\ldots,(x_n)$ with respect to $(P(u,v))$, whereas for $\lambda=1,\ldots,i-1$ the dependences between $(x_\lambda)$ and $(P(u,x_{i+1},\ldots,x_n))$ remain as such, and hence do not further lower the algebraic rank. The algebraic rank has decreased by exactly one unit. If $\frakp$ was of dimension zero, $\frakc$ becomes the unit ideal; the above result therefore remains valid. In the original domain $\bar{\frakS}$, consequently,
\[
 \bar{\fraka}=(\bar{\frakp},v+v_1y_1+\cdots+v_ny_n),
\]
where $v_\lambda$ is put for $-t_{\lambda i}$, is an ideal of the required kind. If, in particular, $P$ contains infinitely many elements, then $v,v_1,\ldots,v_n$ can always be specialized to elements of $P$ in such a way that norm and elementary-divisor form, and hence the dimension of $\bar{\fraka}$, become equal to those of the specialized ideal $\frakb$;\footnote{Compare H.-N., the last paragraph of § 7.} the elementary-divisor form need not, however, remain a prime function. By Theorem V, $\frakb$ nevertheless possesses at least one associated prime ideal of the desired dimension; passing back again to $\bar{\frakS}$, Theorem VI therefore holds there without adjoining any indeterminates.

Theorem II and Theorem VI immediately give the desired formulation according to

\textbf{Theorem VII.} \emph{A prime ideal $\frakp$ has dimension $n-i$ if -- after possible adjunction of indeterminates -- there exists at least one chain of $1+(n-i)+1$ prime ideals}
\[
 \frakp_0=\frakp;\quad \frakp_1,\ldots,\frakp_{n-i};\quad \frakp_{n-i+1},
\]
\emph{each of which is a proper divisor of the immediately preceding one, whereas no such chain with a larger number of members exists. This definition also holds in the original domain, and without introducing any indeterminates, when $P$ contains infinitely many elements.}\footnote{This chain definition of dimension gives, in the case of an algebraic number field, dimension zero for the ordinary prime ideals, dimension $-1$ for the unit ideal, and dimension $1$ for the zero ideal. Indeed the latter also satisfies the definition of prime ideals: in a field a product of non-vanishing quantities is always different from zero. Theorem II remains valid, with this definition of dimension, in algebraic number fields.}

First it is clear that the last member of the chain must be equal to the unit ideal -- which satisfies the definition of a prime ideal -- since otherwise the chain could be lengthened by adjoining $\frako$; for $\frako$ itself Theorem VII gives dimension $-1$, in agreement with § 1, 8. Now let $\frakp$ be different from $\frako$ and of dimension $n-i$. Then the number of members of the chain can be at most $1+(n-i)+1$, since by Theorem II two prime ideals of the same dimension cannot occur. Conversely, by Theorem VI, after possible adjunction of new indeterminates, there exists at least one chain of this length, since the dimension has decreased by exactly one unit at every member of the chain. If one takes Theorem VII as the definition of dimension -- a definition which plainly can be stated in exactly the same way in the original domain, and for which, by Theorem VI, the introduction of indeterminates becomes unnecessary when $P$ contains infinitely many elements -- then the definition of the highest dimension of an arbitrary ideal is again given by the second part of Theorem V.

\subsection*{§ 5. Classification of the Fundamental Ideals and of Their Components in General Ideal Theory. Zeroes of Ideals.}

The question is the characterization of the fundamental ideals and their components (§ 1, 7.) by isolated components of the decomposition (§ 1, 13.) in the sense of general ideal theory. This characterization will show the parallelism between the decomposition theorems for norms and those of general ideal theory, and will make it possible to state, for arbitrary ideals, the theory of zeroes given in § 3 for prime ideals and primary ideals. For the fundamental ideals one has

\textbf{Theorem VIII.} \emph{The fundamental ideals $\frakg_i$ of $i$-th stage are the isolated components of the decomposition that are uniquely determined by the totality of the associated prime ideals of dimensions $n-1,n-2,\ldots,n-i$. If all associated prime ideals have the same dimension, then the ideal also has this dimension; that is, only one factor $R^{(i)}$ of the norm occurs. In general the norm has as many factors $R^{(i)},R^{(i+\sigma)},R^{(i+\tau)},\ldots$ different from unity as different dimension numbers occur among the associated prime ideals.}

Precede the proof by

\textbf{Lemma VI.} \emph{If an ideal is represented as a least common multiple, $\frakm=[\fraka_1,\ldots,\fraka_\alpha]$, then its fundamental ideal $\frakg_i$ is the least common multiple of the fundamental ideals $\frakh_i$ of $i$-th stage of the $\fraka$.}

For from
\[
 b^{(i+1)}f\equiv0(\frakm),\qquad b^{(i+1)}\ne0
\]
one obtains
\[
 b^{(i+1)}f\equiv0(\fraka_\lambda),\qquad b^{(i+1)}\ne0;
\]
therefore $\frakg_i$ is divisible by $[\frakh_{i1},\ldots,\frakh_{i\alpha}]$. Conversely, from
\[
 c_\lambda^{(i+1)}f\equiv0(\fraka_\lambda),\qquad c_\lambda^{(i+1)}\ne0
\]
one obtains
\[
 c_1^{(i+1)}\cdots c_\alpha^{(i+1)}f\equiv0(\frakm),\qquad
 c_1^{(i+1)}\cdots c_\alpha^{(i+1)}\ne0,
\]
and hence the converse divisibility, proving the lemma.

Now let the associated prime ideals of $\frakm$ be ordered by dimension (some dimensions may of course be absent, $j_\lambda=0$):
\[
 \frakp_{1,1},\ldots,\frakp_{1,j_1};\quad
 \frakp_{2,1},\ldots,\frakp_{2,j_2};\quad\ldots;\quad
 \frakp_{n,1},\ldots,\frakp_{n,j_n},
\]
where in general $\frakp_{\lambda,\kappa}$ denotes a prime ideal of dimension $n-\lambda$; let $\frakq_{\lambda,\kappa}$ be the corresponding primary ideal occurring in a fixed given decomposition of $\frakm$. Then by definition the fundamental ideal of $i$-th stage is equal to the unit ideal for all $\frakq_{\lambda,\kappa}$ for which $\lambda>i$, whereas for $\lambda\le i$ it is, by Theorem I, equal to $\frakq_{\lambda,\kappa}$. Therefore, by Lemma VI:
{\let\veqno\leqno\begin{equation}
 \tag{5}
 \frakg_i=[\frakq_{1,1},\ldots,\frakq_{1,j_1};\ldots;\frakq_{i,1},\ldots,\frakq_{i,j_i}],
\end{equation}}
and this representation recognizes $\frakg_i$ as an isolated component of the decomposition, since none of the prime ideals belonging to $\frakg_i$ can be contained in one of the remaining prime ideals, all of which have lower dimension.

If, in particular, only associated prime ideals of dimension $n-i$ occur, then by (5) one has $\frakg_0,\ldots,\frakg_{i-1}$ equal to the unit ideal, while $\frakg_i=\frakg_{i+1}=\cdots=\frakm$; hence $R_\frakm=R^{(i)}$. If, however, prime ideals of different dimensions occur, say $j_i,j_{i+\sigma},j_{i+\tau},\ldots$ are non-zero, then by (5):
\[
 \frakg_0\ne\frakg_i\ne\frakg_{i+\sigma}\ne\frakg_{i+\tau}\ne\cdots,
\]
whereas
\[
 \frakg_0=\frakg_1=\cdots=\frakg_{i-1}=\frako,
 \qquad
 \frakg_i=\frakg_{i+1}=\cdots=\frakg_{i+\sigma-1},\ldots,
\]
therefore the factors $R^{(i)},R^{(i+\sigma)},R^{(i+\tau)}$, and only these, are different from unity.\footnote{Taking Lemma II into account, Theorem VIII again implies that the fundamental ideals -- as least common multiples of transformed ideals -- become transformed ideals. Since Theorem VIII rests only on theorems from H.-N. where this fact is not used, this gives at the same time a new proof which reveals the internal reason. In H.-N. the theorem is used only in the theory of zeroes, which is newly developed here.}

By means of Theorem VIII, Theorem I sharpens to

\textbf{Theorem IX.} \emph{An ideal is primary if and only if its elementary-divisor form -- and therefore its norm -- is a primary function.}

That the condition is fulfilled for primary ideals was shown in Theorem I. Conversely, let the elementary-divisor form of $\frakm$ be given by $Q^{(i)}=(P^{(i)})^\rho$, where $P^{(i)}$ denotes a prime function. By Theorem VIII all prime ideals of $\frakm$ then have the same dimension $n-i$. From $\frakm=[\frakq_1,\ldots,\frakq_a]$ one obtains:
\[
 Q^{(i)}\equiv0(\frakq_\lambda);\quad\hbox{hence}\quad
 (P^{(i)})^\rho\equiv0(\frakp_\lambda)\quad\hbox{and therefore}\quad
 P^{(i)}\equiv0(\frakp_\lambda).
\]
Since all $\frakp_\lambda$ have dimension $n-i$ and $P^{(i)}$ is a prime function, $P^{(i)}$ becomes the elementary-divisor form of each $\frakp_\lambda$. Thus, by the corollary to Theorem IV, all $\frakp_\lambda$ coincide. Hence, since greatest primary components are involved, only one $\frakp$ can occur; $\frakm$ is primary, the special prime case for $\rho=1$ not being excluded. Theorems VIII and IX lead to the classification of the components of the fundamental ideals as follows:

\textbf{Theorem X.} \emph{The component $\frakr_\lambda$ of the fundamental ideal $\frakg_i$ corresponding to the greatest primary factor $Q_\lambda^{(i)}=(P_\lambda^{(i)})^{\rho_\lambda}$ of $R^{(i)}$ is given by the isolated component of the decomposition which is uniquely determined by the associated prime ideals of dimensions $n-1,\ldots,n-i+1$ and by that prime ideal of dimension $n-i$ whose elementary-divisor form becomes equal to $P_\lambda^{(i)}$. Associated prime ideals and greatest primary factors of the norm correspond one-to-one.}

By § 1, 7. the $\frakr_\lambda$ agree with their fundamental ideal of $i$-th stage and possess $\frakg_{i-1}$ as fundamental ideal of $(i-1)$-st stage. Hence, by Theorem VIII respectively (5), they admit a shortest representation:
\[
 \frakr_\lambda=[\frakg_{i-1},\frakq_{\lambda,1},\ldots,\frakq_{\lambda,t_\lambda}],
\]
where all $\frakq$ have dimension $n-i$ and belong to distinct prime ideals. By the definition of $Q_\lambda^{(i)}$ -- taking into account that $\frakr_\lambda$ agrees with its fundamental ideal of $i$-th stage -- one has:
\[
 Q_\lambda^{(i)}\frakg_{i-1}\equiv0(\frakq_{\lambda,\kappa}),
 \quad(\kappa=1,\ldots,t_\lambda)
\]
and consequently
\[
\begin{aligned}
 (Q_\lambda^{(i)})^{\sigma_\lambda}&\equiv0(\frakq_{\lambda,\kappa});
 \quad\hbox{hence}\quad
 (P_\lambda^{(i)})^{\rho_\lambda\sigma_\lambda}\equiv0(\frakp_{\lambda,\kappa}),\\
 &\hbox{and therefore}\quad P_\lambda^{(i)}\equiv0(\frakp_{\lambda,\kappa}),
\end{aligned}
\]
from which, as in the proof of Theorem IX, it follows that in the representation of $\frakr_\lambda$ only one $\frakq_\lambda$ and associated $\frakp_\lambda$ of dimension $n-i$ can occur, and that the elementary-divisor form of $\frakp_\lambda$ is given by $P_\lambda^{(i)}$. It remains to prove that $\frakp_\lambda$ is an associated prime ideal of $\frakm$ and that $\frakq_\lambda$ actually occurs in at least one representation of $\frakm$ by greatest primary components. Then $\frakr_\lambda$ is also recognized as the isolated component of the decomposition, since no prime ideal can be absorbed in one of equal or lower dimension, namely as the component determined by $\frakp_\lambda$ and by all associated prime ideals of higher dimension. For this proof it suffices to show that $\frakq_\lambda$ occurs in a shortest representation of $\frakg_i$, since such a representation can always, by Theorem VIII, be completed to a shortest representation of $\frakm$. Thus it is enough to show that the representation
\[
 \frakg_i=[\frakr_1,\ldots,\frakr_\alpha]=[\frakg_{i-1},\frakq_1,\ldots,\frakq_\alpha]
\]
becomes a shortest one. If, for instance, $\frakq_\lambda$ were absorbed by its complement, so that $\frakg_{i-1}\frakq_1\cdots\frakq_{\lambda-1}\frakq_{\lambda+1}\cdots\frakq_\alpha$ were divisible by $\frakq_\lambda$, then, because $\frakg_{i-1}\not\equiv0(\frakq_\lambda)$, it would follow that a power of some $\frakq_\mu$ is divisible by $\frakq_\lambda$ for $\mu\ne\lambda$. Thus the associated prime ideals $\frakp_\mu$ and $\frakp_\lambda$, both of the same dimension, would have to coincide by Theorem II; but this is impossible, since their elementary-divisor forms $P_\mu^{(i)}$ and $P_\lambda^{(i)}$ are distinct by definition. The representation of $\frakg_i$, to which by Theorem VIII all associated prime ideals of dimension $n-i$ correspond, further shows that to every such prime ideal there really corresponds a greatest primary factor of the norm. The one-to-one correspondence is therefore proved.

Theorem X immediately gives the transfer of Theorem IV according to

\textbf{Theorem XI.} \emph{The individual greatest primary factors $Q_\lambda^{(i)}=(P_\lambda^{(i)})^{\rho_\lambda}$ of the norm admit a product representation:}
\[
 Q_\lambda^{(i)}=\prod (x_i-t_{1i}\bar y_{1\nu}-\cdots-t_{ni}\bar y_{n\nu})^{\delta_\lambda\rho_\lambda}
 =\prod (t_{1i}(y_1-\bar y_{1\nu})+\cdots+t_{ni}(y_n-\bar y_{n\nu}))^{\delta_\lambda\rho_\lambda},
\]
\emph{where $\bar y_{1\nu},\ldots,\bar y_{n\nu}$ represents in each case an associated system of zeroes of the associated prime ideal $\frakp_\lambda$, in the original indeterminates $y$, independent of $t_{1i},\ldots,t_{ni}$, and $\delta_\lambda$ has the value one or $p^{f_\lambda}$. Thus the decomposition of the norm into linear factors is completely given. The degree of $Q_\lambda^{(i)}$ gives the degree, over the residue-class subfield $(P(u,x_{i+1},\ldots,x_n))$ obtained by quotient formation, of the subring of residue classes, represented by residue classes from $\frakg_{i-1}$, modulo the isolated component $\frakr_\lambda$; for primary ideals this is therefore the ring of all residue classes. When new indeterminates $v$ are adjoined, one further obtains the decomposition:}
\[
 Q_\lambda^{(i)}(u,v)=\prod (z-v_1\bar x_{1\nu}-\cdots-v_i\bar x_{i\nu})^{\delta_\lambda\rho_\lambda}.
\]

Since $P_\lambda^{(i)}$ has been recognized by Theorem X as the elementary-divisor form of $\frakp_\lambda$, this is in fact only the substitution of the decomposition of Theorem IV. The remark on the degree is only another formulation of what was said in § 1, 5. and 7.\footnote{Likewise the formulation of multiplicity mentioned in note 2) of H.-N. is an immediate consequence of Theorem X. For, after adjoining $x_{i+1},\ldots,x_n$, the residue-class system $\frakg_{i-1}\mid\frakr_\lambda$ is isomorphic to $\frakt_\lambda\mid\frakr_\lambda$, where $\frakt_\lambda=[\frakg_{i-1},\frakr_1,\ldots,\frakr_{\lambda-1},\frakr_{\lambda+1},\ldots,\frakr_a]$ is put, as follows from passing back to modules of linear forms (H.-N., Theorem V) in view of the relative primality of the factors $Q_\mu^{(i)}$. Directly it is shown that $\frakt_\lambda\mid\frakr_\lambda$ is isomorphic to $\frakt_\lambda\mid\frakq_\lambda$; here $\frakt_\lambda$ arises from the complement of $\frakq_\lambda$ after adjoining $x_{i+1},\ldots,x_n$.}

\subsection*{§ 6. Characterization of Prime Ideals and Proper Primary Ideals by Elementary-Divisor Form and Norm.}

In Theorem IX primary ideals were characterized by elementary-divisor form and norm, but only in such a way that the prime case was to be regarded as a special case of primary. We now need to separate prime ideals from proper primary ideals; the considerations will run parallel to those of § 3. Conditions that are either necessary or sufficient are easily stated according to

\textbf{Theorem XII.} \emph{A necessary condition for a prime ideal is that its elementary-divisor form be a prime function; a sufficient condition is that its norm be a prime function. In other words, the elementary-divisor form of a prime ideal is always a prime function, and the norm of a proper primary ideal is always a proper primary function.}

That the elementary-divisor form of a prime ideal is a prime function was already shown in Theorem I. Suppose now that the norm $R_\frakq$ is equal to a prime function. It is to be shown that, under this hypothesis, every proper divisor $\fraka$ of $\frakq$ has lower highest dimension; then $\frakq$ is recognized as a prime ideal by Theorem II. Indeed, by § 1, 6. the first factor of the norm $R_\fraka$ that differs from the corresponding one of $R_\frakq$ is a proper divisor. Since $R_\frakq=P^{(i)}$, the factor $R^{(i)}$ of $R_\fraka$ must therefore be equal to a proper divisor of $P^{(i)}$, hence to unity; $\fraka$ has lower highest dimension.

A second simple proof rests on the following lemma, on which the next theorem also rests:

\textbf{Lemma VII.} \emph{The system of residue classes represented by polynomials $H^{(i)}$ modulo the elementary-divisor form $E^{(i)}$ of a primary ideal -- more generally, of an ideal that has only associated prime ideals of dimension $n-i$ -- is isomorphic to a subsystem of the residue classes modulo this ideal. The degree of the elementary-divisor form gives the degree of this residue-class system with respect to the quotient field $(P(u,x_{i+1},\ldots,x_n))$.}

Indeed, one defines a correspondence by assigning to one another the classes represented by the same polynomials $H^{(i)}$; sums and products correspond to sums and products. The assignment is also one-to-one. For all polynomials $H^{(i)}$ belonging to the zero class modulo the ideal are, by definition, divisible by $E^{(i)}$; hence the zero class corresponds to the zero class modulo $E^{(i)}$, and conversely as well.

Now let the norm of an ideal be a prime function $P^{(i)}$, hence identical with the elementary-divisor form. The degrees of the residue-class systems modulo the elementary-divisor form $P^{(i)}$ and modulo $\frakq$ therefore agree with respect to $(P(u,x_{i+1},\ldots,x_n))$; the subsystem isomorphic to the residue-class system modulo $P^{(i)}$ therefore exhausts the residue classes modulo $\frakq$, and since this system can have no zero divisors, $\frakq$ is a prime ideal.

In general Theorem XII cannot be sharpened to necessary and sufficient conditions, as will be shown below by examples. This is possible, however, when the coefficient domain $P$ is a perfect field (§ 1, 14.), according to

\textbf{Theorem XIII.} \emph{If the coefficient domain $P$ is a perfect field, then norm and elementary-divisor form of a prime ideal become prime functions and consequently identical; norm and elementary-divisor form of a proper primary ideal become proper primary functions. Thus, over a perfect field, the condition that the elementary-divisor form be a prime function is necessary and sufficient for a prime ideal; in the condition the norm may also stand in place of the elementary-divisor form.}

Taking Theorem XII into account, Theorem XIII will be proved once it is shown that over a perfect field the ideal becomes prime as soon as its elementary-divisor form is a prime function, and that then the norm is equal to this prime function. In Lemma V, § 3, it was shown that, over a perfect field, the elementary-divisor form becomes a prime function of the first kind as soon as it is a prime function at all. Therefore let $(\mR)$ denote the residue-class field obtained by adjoining $(x_i)$ to $(P(u,x_{i+1},\ldots,x_n))$, that is, $(P(u,x_i,x_{i+1},\ldots,x_n))$ modulo $P^{(i)}$. Then $(x_i)$ is a zero of the prime function of the first kind $(P^{(i)})$, and hence, as for example Lagrange's formula shows, $(y_1),\ldots,(y_n)$ are contained in $(\mR)$. The subsystem of the residue classes modulo the ideal $\frakq$ under consideration that is isomorphic, by Lemma VII, to $(\mR)$ -- where only denominators from $(P(u,x_{i+1},\ldots,x_n))$ occur -- therefore contains the residue classes $(y_1),\ldots,(y_n)$ and hence exhausts the residue-class system modulo $\frakq$. Since, being isomorphic to $(\mR)$, it can have no zero divisors, $\frakq$ is a prime ideal. The degree of this residue-class field $(\mR)$ of $\frakq$ with respect to $(P(u,x_{i+1},\ldots,x_n))$ is equal to the degree of the norm of $\frakq$; on the other hand, by the isomorphism with $(\mR)$, it is equal to the degree of the elementary-divisor form $P^{(i)}$. Thus norm and elementary-divisor form must agree.

For imperfect coefficient domains Theorem XII still allows the following sharpening:

\textbf{Theorem XIV.} \emph{If the coefficient domain $P$ is an imperfect field of characteristic $p$, then the norm of a prime ideal is the $p^g$-th power of the elementary-divisor form; one has $0\le g\le(i-1)f$, where $f$ denotes the exponent of the elementary-divisor form and $n-i$ the dimension of the prime ideal.}

Theorem XIV amounts to the assertion that the residue-class field $(\mR)$ of $\frakp$ has degree $p^g$ ($g\ge0$) with respect to the subfield $(\mR')=(P(u,x_i,x_{i+1},\ldots,x_n))$ isomorphic to $(\mR)$. This follows directly from the Steinitz theorems stated in § 1, 15., more precisely from

\textbf{Lemma VIII.} \emph{Let $\Lambda$ be a finite extension of the imperfect field $\Omega$ of reduced degree $r$ and exponent $f$; let $\Lambda_0$ be the field of the first kind contained in $\Lambda$, and let $M$ be an intermediate field between $\Lambda_0$ and $\Lambda$. Then the degree of every element of $\Lambda$ with respect to $M$ is a power of $p$, where $p$ denotes the characteristic of $\Omega$.}

For the $p^{f'}$-th power ($f'\le f$) of every element $z$ of $\Lambda$ belongs to $\Lambda_0$; hence $z$ is a zero of a prime function $G(t)=(t-z)^{p^{f'}}$ with coefficients in $\Lambda_0$. Let $H(t)=(t-z)^\delta$ be the prime function in $M$ with zero $z$. Then $H(t)$ also has coefficients in $\Lambda_0(z)$, and therefore the coefficients of $H(t)$ belong to the intersection field of $M$ and $\Lambda_0(z)$, hence to an intermediate field between $\Lambda_0$ and $\Lambda_0(z)$. The degree of $z$ with respect to $M$ is therefore equal to the degree with respect to this intermediate field, and hence, as a divisor of $p^{f'}$, is a power of $p$.

For the proof of Theorem XIV it remains only to observe that the field $(\mR')=(P(u,x_i,x_{i+1},\ldots,x_n))$ must have degree $r\cdot p^f$ with respect to $(P(u,x_{i+1},\ldots,x_n))$, if $r$ denotes the reduced degree and $f$ the exponent of $(\mR)$; and consequently the field of the first kind $(\mR_0)$ contained in $(\mR)$ must be a subfield of $(\mR')$. Since there are elements of degree $r\cdot p^f$ in $(\mR)$ (§ 1, 15.), $(x_i)$ must necessarily be of this degree, since all elements of $(\mR)$ arise by specializing the $t_{\mu\nu}$ in $(x_i)$. The exponent of $(\mR)$ agrees with the exponent of the elementary-divisor form, and $(x_i)^{p^f}$ has degree $r$ and is an element of $(\mR_0)$, hence a primitive element of $(\mR_0)$. Since $(\mR)$ arises by adjoining finitely many elements -- say $(x_1),\ldots,(x_{i-1})$ -- to $(\mR')$, repeated finite application of Lemma VIII gives the desired proof and at the same time the estimate for $g$.\footnote{Compare the parallel considerations in A. Ostrowski, Zur arithmetischen Theorie der algebraischen Größen, Gött. Nachr. 1919, pp. 279--298, especially p. 288, where the corresponding theorem is stated without proof.}

We now add the examples which show that Theorem XII cannot in general be sharpened beyond Theorem XIV. They concern a prime ideal whose norm is equal to the square of the elementary-divisor form ($g>0$), and proper primary ideals whose elementary-divisor forms are prime functions. The latter example also shows that no analogue of Theorem XIV exists here, but that the norm can become an arbitrary power of the elementary-divisor form.

Let the coefficient domain $P$ arise by adjoining two indeterminates $\lambda,\mu$ to the field of residue classes modulo $2$, and therefore let $P$ be an imperfect field of characteristic two. Consider the ideal
\[
 \bar{\frakp}=(y_1^2+\lambda,\,y_2^2+\mu),
\]
which must be a prime ideal, since the residue-class system $(\bR)$ is isomorphic to the field $P(\sqrt\lambda,\sqrt\mu)$. The field $(\bR)$ has degree four with respect to $(P)$, reduced degree $r=1$, and exponent $f=2$. Therefore the elementary-divisor form $P^{(2)}$ is of the second degree, and the norm is of the fourth degree and consequently the square of $P^{(2)}$, as can also be checked directly from the module representation (3), § 1, 5. One obtains
\[
 P^{(2)}=(u_{11}u_{22}+u_{12}u_{21})^2x_2^2+(u_{11}^2\mu+u_{21}^2\lambda)
\]
or, after a suitable multiplication:
\[
 x_2^2+(t_{12}^2\lambda+t_{22}^2\mu).
\]

Let $P$ now arise by adjoining the one indeterminate $\lambda$ to the field of residue classes modulo $2$, and take as basis the ideal
\[
 \bar{\frakq}=(y_1^2+\lambda,\,y_2^2+\lambda)=(y_1^2+\lambda,(y_1+y_2)^2)
\]
As elementary-divisor form, by specializing $\mu=\lambda$ in the one just given, one obtains:
\[
 P^{(2)}=x_2^2+\lambda(t_{12}^2+t_{22}^2),
\]
hence a prime function, since $t_{12}=(0)$, $t_{22}=(1)$ leads to the prime function $x_2^2+\lambda$. But $\bar{\frakq}$ is proper primary, since it contains $(y_1+y_2)^2$ but not $(y_1+y_2)$. The norm must be equal to the square, hence to the $p$-th power, of $P^{(2)}$, since the number of residue classes of $\bar{\frakq}$ linearly independent over $P$ is four.

That this last analogue of Theorem XIV is not always fulfilled is shown by the following example. Let $P$ arise by adjoining the indeterminate $\lambda$ to the residue-class field modulo $3$, and take as basis
\[
 \bar{\frakq}=(y_1^3+\lambda,(y_1-y_2)^2)
\]
so that, as above, $\bar{\frakq}$ is a proper primary ideal. The elementary-divisor form becomes -- taking into account that $y_2^3+\lambda$ is also divisible by $\bar{\frakq}$ --
\[
 P^{(2)}=x_2^3+\lambda(t_{12}+t_{22})^3,
\]
a prime function; but the norm is equal to the square of $P^{(2)}$, since there are six linearly independent residue classes modulo $\bar{\frakq}$. Thus the $p^g$-th power of the elementary-divisor form does not occur here.

The first example shows that, over an imperfect field as coefficient domain, just as in algebraic number fields, prime ideals of degree higher than the first may occur, where $p^g$ is to be called the degree of the prime ideal. The further examples are to be understood as analogues of ramification ideals: a prime number may be divisible by a higher power of a prime ideal, by a primary ideal. For, as shown in note 10), the elementary-divisor form corresponds to the least integral rational number divisible by an ideal in an algebraic number field.

\subsection*{§ 7. Absolute Prime Ideals.}

A prime ideal with coefficients in $P$ is called an \emph{absolute prime ideal} if it remains a prime ideal in the algebraically closed field $A$ to which $P$ can be extended.\footnote{An algebraically closed field is, as is well known, a field in which every polynomial in one indeterminate decomposes into linear factors. For the existence and essential uniqueness of the algebraically closed field belonging to an arbitrary field, see Steinitz.} Correspondingly, a prime function $P$ with coefficients in $P$ is called an \emph{absolute prime function} -- absolutely irreducible polynomial -- if $P$ remains a prime function in $A$. The characterization of absolute prime ideals is given by

\textbf{Theorem XV.} \emph{A necessary and sufficient condition for an absolute prime ideal is that its elementary-divisor form -- which may be assumed integral and primitive in the $u$ -- become an absolute prime function as a polynomial in the $x$ and $u$. The elementary-divisor form becomes identical with the norm, so that the condition may also be stated for the norm.}

For the proof, observe first that in general norm and elementary-divisor form are preserved under algebraic extension of the coefficient domain. Indeed the individual factors $R^{(i)}$ respectively $E^{(i)}$ are defined as greatest common divisors in the polynomial sense, a property preserved by algebraic extension of the coefficient domain. Thus $P^{(i)}$ remains the elementary-divisor form of $\frakp$ when passing from $P$ to $A$, hence when passing to a perfect field as coefficient domain, since every algebraically closed field is perfect, as the definition immediately shows. By Theorem XIII, a necessary and sufficient condition for $\frakp$ to be an absolute prime ideal is that $P^{(i)}$ become a prime function with respect to $A(u)$; that is, an absolute prime function as a polynomial in the $x$ and $u$, if $P^{(i)}$ is assumed integral and primitive in the $u$. At the same time, by Theorem XIII, $P^{(i)}$ becomes identical with the norm of $\frakp$ as soon as $A$ is taken as coefficient domain. Hence, by what was just noted, the same holds over $P$. This proves Theorem XV.

For absolute prime functions $P$ with coefficients in a (finite) algebraic number field $\mathfrak K$, one now has the theorem\footnote{A. Ostrowski, loc. cit. (note 21)); lemma p. 296. -- A simpler proof is in E. Noether, Ein algebraisches Kriterium für absolute Irreduzibilität, Math. Ann. 85 (1922), pp. 26--33, no. 7.} that $P$ remains a prime function, more precisely an absolute prime function, modulo every prime ideal of $\mathfrak K$, with at most finitely many prime ideals of $\mathfrak K$ excepted. The transfer of this theorem to absolute prime ideals rests on

\textbf{Theorem XVI.} \emph{If, as coefficient domain, one takes, in place of a field, the ring $\mathfrak o^*$ of all algebraic integers of a (finite) algebraic number field $\mathfrak K$, then norm and elementary-divisor form of a polynomial ideal are preserved as norm and elementary-divisor form modulo every prime ideal of $\mathfrak K$, with at most finitely many prime ideals of $\mathfrak K$ excepted.}

For the proof the polynomial domain on which everything is based is to be modified relative to § 1, 1. Let $\bar{\frakS}$ consist of all polynomials in $y_1,\ldots,y_n$ with coefficients in $\mathfrak o^*$, hence algebraic integers of $\mathfrak K$; let the coefficient domain of $\frakS$ be $\mathfrak o^*[u]$, that is, all polynomials in $u$ with coefficients in $\mathfrak o^*$. If $\bar{\frakm}$ is an ideal in $\bar{\frakS}$, it passes by (1) into a transformed ideal $\frakm$ in $\frakS$. It is to be shown that only finitely many polynomials from $\mathfrak o^*[u]$ occur in the denominators in forming the norm; then Theorem XVI follows directly.

First note that the module representation (3) in § 1, 5. for the formation of the individual norm amounts to reducing powers of $x_{i-1}$ modulo a polynomial regular in $x_{i-1}$, $C^{(i-1)}=U_{i-1}(u)x_{i-1}^{k}+\hbox{lower terms}$\footnote{Compare H.-N. § 4.}; here all coefficients of $C^{(i-1)}$, in particular $U_{i-1}$, may be assumed to be elements of $\mathfrak o^*[u]$. Since this is a matter of successive reduction of $x_1,\ldots,x_{i-1}$, the module representation integral in the $x$ is also integral in $\mathfrak o^*[u]$ except for power products $U_1^{\lambda_1}\cdots U_{i-1}^{\lambda_{i-1}}$ in the denominator. The product $U_1U_2\cdots U_n$, considered as a polynomial in the $u$, defines by its coefficients an ideal $\mathfrak r^*$ of $\mathfrak K$, which is therefore divisible by only finitely many prime ideals of $\mathfrak K$. If $\frakp^*$ is chosen different from these finitely many, and if the residue-class field modulo $\frakp^*$ is used as coefficient domain, the original module representation is preserved by replacing each number of $\mathfrak o^*$ by its residue class modulo $\frakp^*$.

Furthermore, since norm and elementary-divisor form are determined only up to factors from $\mathfrak o^*[u]$, they may also be assumed divisible by $\frakm$ with respect to $\mathfrak o^*[u]$. Suppose, under this convention, that for instance $R^{(i)}=V_i(u)x_i^{l}+\hbox{lower terms}$, so that, in the divisibility of all $\rho$-rowed determinants of the module $\mathfrak M_{i-1}^*$ of rank $\rho$ determined by $C^{(i-1)}$, only powers of $V_i$ occur in the denominator besides the $U_1^{\lambda_1}\cdots U_{i-1}^{\lambda_{i-1}}$. Choose $\frakp^*$ further distinct from the finitely many prime ideals of $\mathfrak K$ that occur in the ideal determined by $V_1V_2\cdots V_n$. Then $R^{(i)}$ remains a common divisor of these determinants when the residue-class field modulo $\frakp^*$ is taken as coefficient domain, and the rank of $\mathfrak M_{i-1}^*$ is preserved, since $C^{(i)}$ was a $\rho$-rowed determinant from $\mathfrak M_{i-1}^*$.

Finally $\frakp^*$ is to be chosen so that $R^{(i)}$ remains the greatest common divisor in each case. Let $D^{(i)}$ run through all $\rho$-rowed determinants from $\mathfrak M_{i-1}^*$, so that by hypothesis $U_1^{\lambda_1}\cdots U_{i-1}^{\lambda_{i-1}}V_i^{\sigma}D^{(i)}=R^{(i)}T^{(i)}$, with the $T^{(i)}$ having no common divisor containing $x$ in the polynomial sense. Consequently the ideal derived from all $T^{(i)}$ contains a polynomial $G^{(i+1)}=W_i(u)x_{i+1}^{m}+\hbox{lower terms},\ W_i\ne0$; that $G^{(i+1)}$ may be assumed regular in $x_{i+1}$ follows from composing the transformation (1) with one of $x_{i+1},\ldots,x_n$ having new indeterminates as coefficients. If, then, $\frakp^*$ is also chosen distinct from the finitely many prime ideals contained in the ideal determined by $W_1W_2\cdots W_n$, then when the residue-class field modulo $\frakp^*$ is used as coefficient domain, $R_\frakm$ is in fact preserved as norm. Quite correspondingly $\frakp^*$ can be chosen still further so that $E_\frakm$ also remains the elementary-divisor form. This proves Theorem XVI.

From Theorems XV and XVI, as a generalization of the theorem on absolute prime functions and using that theorem, one obtains

\textbf{Theorem XVII.} \emph{If $\frakp$ is an absolute prime ideal with algebraic integers from a (finite) algebraic number field $\mathfrak K$ as coefficients, then $\frakp$ remains a prime ideal, more precisely an absolute prime ideal, modulo every prime ideal of $\mathfrak K$, with at most finitely many prime ideals of $\mathfrak K$ excepted.}

For the proof, let the norm $R_\frakp$ of $\frakp$ be chosen, as in Theorem XVI, so that it is divisible by $\frakp$ also with respect to $\mathfrak o^*[u]$. Then
$R_\frakp=T(u)P^{(i)}(x)$, where $P^{(i)}(x)$ may be assumed integral and primitive in the $u$. By Theorem XV, $P^{(i)}$, regarded as a polynomial in $u$ and $x$ with coefficients in $\mathfrak K$, is an absolute prime function. By the theorem on absolute prime functions it retains this property modulo every prime ideal of $\mathfrak K$, with at most finitely many prime ideals of $\mathfrak K$ excepted. Thus choose $\frakp^*$ different from these finitely many prime ideals and, by Theorem XVI, also different from finitely many further prime ideals of $\mathfrak K$. Then $P^{(i)}$ -- its coefficients from $\mathfrak o^*$ being replaced by the corresponding residue classes modulo $\frakp^*$ -- becomes the norm of $\frakp$ when the residue-class field modulo $\frakp^*$ is taken as coefficient domain $P$, and this norm is an absolute prime function. Taking $P$ as coefficient domain, $\frakp$ therefore remains an absolute prime ideal by Theorem XV. This proves Theorem XVII.

\begin{center}
Göttingen, May 8, 1923.\\[1em]
(Received March 9, 1923.)
\end{center}



\clearpage

\end{document}
